\chapter{CY-C: Six Routes to \texorpdfstring{$G(K3 \times E)$}{G(K3 x E)} and Their Pairwise Comparison}
\label{ch:cy-c-six-routes-convergence}

\providecommand{\CoHA}{\mathrm{CoHA}}
\providecommand{\Coh}{\mathrm{Coh}}
\providecommand{\HH}{\mathrm{HH}}
\providecommand{\Sym}{\mathrm{Sym}}
\providecommand{\Hilb}{\mathrm{Hilb}}
\providecommand{\Eone}{E_{1}}
\providecommand{\Etwo}{E_{2}}
\providecommand{\LNar}{\Lambda_{\mathrm{Nar}}}
\providecommand{\rhoR}{\rho}
\providecommand{\kch}{\kappa_{\mathrm{ch}}}

Conjecture CY-C is the last open conjecture in the CY-A/B/C/D hierarchy: CY-A\textsubscript{3} is proved at the $\infty$-categorical level, CY-B is proved in its $d$-dependent form, and CY-D is proved in its dimension-stratified form.  What remains is the assertion that the six independent constructions of a quantum vertex chiral group $G(K3 \times E)$ converge on a single isomorphism class.

Only one of the six routes is an application of the CY-to-chiral correspondence $\Phi_3 = \SpCh_{\Sigma_2, C} \circ \PhiFA_3$ (Theorem~\ref{thm:phi-two-stage-factorisation}; Definition~\ref{def:phi-fa-and-sp}); the remaining five are independent constructions whose outputs carry their own modular characteristics, their own $R$-matrices, and their own automorphic data.  In the two-stage framework, CY-C is the conjecture that the outputs of the five other routes R\textsubscript{2}--R\textsubscript{6} admit comparison maps to $\Eone$-chiral shadows that should be representable by distinct $(\Sigma_2, C)$-specialisations of the canonical $E_3$-holomorphic factorisation algebra $\PhiFA_3(D^b\Coh(K3 \times E))$ on $K3 \times E$ (Remark~\ref{rem:multiple-e1-specialisations}).  The comparison maps, their compatibility, and the resulting pairwise isomorphisms are the \emph{content} of CY-C, not a derivation from functoriality.

\section{Setup: the six constructions}
\label{sec:cy-c-six-setup}

Fix a K3 surface $S$ with N\'eron--Severi lattice of signature $(1, \rho(S) - 1)$, an elliptic curve $E$, and the Calabi--Yau threefold $X = S \times E$. The six routes extract a chiral algebra (or chiral group) from this geometry through six different machines.

\begin{definition}[The six routes]
\label{def:cy-c-six-routes}
\ClaimStatusDefinitional
A \emph{route} for $G(X)$ is a pair (construction, output) producing a chiral algebra $A^{R_i}_X$ and, where applicable, a quantum vertex chiral group $G^{R_i}(X)$. The six canonical routes are:
\begin{enumerate}[label=\textbf{R\arabic*.}]
  \item \textbf{CY-to-chiral correspondence $\Phi_3$.} $\Phi_3 = \SpCh_{\Sigma_2, C} \circ \PhiFA_3 \colon D^b(\Coh(X)) \longrightarrow \mathrm{ChirAlg}^{\Eone}_X$ (Theorem~\ref{thm:phi-two-stage-factorisation}), applied to $D^b(\Coh(S \times E))$, producing $A_X^{R_1} := \Phi_3(D^b(\Coh(X))) = \SpCh_{\Sigma_2, C}(\PhiFA_3(D^b(\Coh(X))))$.  Here $\PhiFA_3(D^b(\Coh(X))) \in \EdHolFA(X)$ is the canonical $E_3$-holomorphic factorisation algebra on $X$, and $\SpCh_{\Sigma_2, C}$ is the stage-$2$ specialisation along a $2$-cycle $\Sigma_2 \subset X$ and reference curve $C$.  \emph{Status of output}: $A_X^{R_1}$ exists as chiral algebra (CY-A\textsubscript{3} proved, Theorem~\ref{thm:cy-to-chiral-d3}); its quantum vertex chiral group structure is conjectural.
  \item \textbf{Borcherds lift.} The Borcherds multiplicative lift of the K3 elliptic genus $2\phi_{0,1}$ produces the Siegel cusp form $\Phi_{10}$ (or $\Delta_5$, its square root up to a constant). The root multiplicities from the denominator identity (Theorem~\ref{thm:k3e-denominator}) define the generalized Borcherds--Kac--Moody superalgebra $\mathfrak{g}_{\Delta_5}$. Output: $A_X^{R_2} := U(\mathfrak{g}_{\Delta_5})$-modules on the natural Mukai lattice completion.
  \item \textbf{Lattice VOA.} The Mukai lattice $H^*(S, \Z) \oplus H^*(E, \Z)$ carries a natural even unimodular quadratic form (Mukai pairing on $S$, standard pairing on $E$), producing a lattice vertex operator algebra $V_{\Lambda_{\mathrm{Muk}}(X)}$. Output: $A_X^{R_3} := V_{\Lambda_{\mathrm{Muk}}(X)}$.
  \item \textbf{Kummer construction.} The Kummer resolution $\Kum(X) := \widetilde{(S \times E)/\iota}$ for $\iota$ the standard involution produces, after orbifold completion, a chiral algebra $A_X^{R_4}$ as the invariant subalgebra of a twisted lattice VOA. The construction is modelled on the classical Kummer K3 but lifted to the threefold via the fibration $S \times E \to E/\iota$.
  \item \textbf{Sigma model.} The $\mathcal{N} = (2,2)$ superconformal sigma model on $X$ admits a topological half-twist producing the holomorphic stress tensor and chiral currents. Output: $A_X^{R_5}$ is the chiral algebra of the holomorphic half-twist on $X$.
  \item \textbf{BLLPR holomorphic dependence.} Beem--Lemos--Liendo--Peelaers--Rastelli (arXiv:1312.5344) associate to a $4$d $\mathcal{N} = 2$ theory a $2$d chiral algebra through the Schur limit. The CY-3 counterpart, discussed in BLLPR-type constructions for class-$S$ theories with K3-fibered geometries (see also arXiv:1710.03254 for the holomorphic sigma model), produces $A_X^{R_6}$ from the holomorphic-twist of a $6$d $(2,0)$ theory compactified on a Riemann surface with K3 structure.
\end{enumerate}
Only route R\textsubscript{1} is an application of $\Phi_3$. Routes R\textsubscript{2}--R\textsubscript{6} are independent constructions.
\end{definition}

\begin{remark}[Why the six routes do not reduce to one]
\label{rem:six-routes-irreducible}
The six constructions use genuinely different mathematical data. R\textsubscript{1} consumes the bounded derived category of coherent sheaves; R\textsubscript{2} consumes a weak Jacobi form of weight $0$ and index $1$; R\textsubscript{3} consumes a rank-$24$ even unimodular lattice with a distinguished signature decomposition; R\textsubscript{4} consumes a finite-group action with specified fixed loci; R\textsubscript{5} consumes a Ricci-flat K\"ahler metric (equivalently, a topological twist of a superconformal theory); R\textsubscript{6} consumes a $6$d superconformal theory compactified on a Riemann surface. None of these inputs refines the others. Convergence of the six outputs is the \emph{content} of CY-C.
\end{remark}

\begin{remark}[The six routes in the two-stage framework]
\label{rem:six-routes-two-stage}
Theorem~\ref{thm:phi-two-stage-factorisation} factors $\Phi_3 = \SpCh_{\Sigma_2, C} \circ \PhiFA_3$: stage~$1$ produces the canonical $E_3$-hFA $\PhiFA_3(D^b(\Coh(X))) \in \EdHolFA(X)$ pinned up to contractible choice by the CY datum, and stage~$2$ specialises along a choice of $(\Sigma_2, C)$ to an $\Eone$-chiral algebra on $C$ (Proposition~\ref{prop:native-en-level}). Only route R\textsubscript{1} is a canonical instance of this two-stage pipeline. Routes R\textsubscript{2}--R\textsubscript{6} are independent constructions; CY-C predicts that, after the missing comparison maps are built, their $\Eone$-chiral shadows correspond to alternative $(\Sigma_2, C)$-specialisations of the same canonical $\PhiFA_3(D^b(\Coh(X)))$ through Remark~\ref{rem:multiple-e1-specialisations}: the Borcherds lift (R\textsubscript{2}) to the K3-fibre / $E$ specialisation, the lattice VOA (R\textsubscript{3}) to the Mukai-lattice Nakajima specialisation, the Kummer orbifold (R\textsubscript{4}) to a $(\Z/2)$-twisted specialisation, the sigma-model half-twist (R\textsubscript{5}) to the Ricci-flat metric specialisation, and the BLLPR chiral algebra (R\textsubscript{6}) to the Schur-limit specialisation of the $6$d $(2,0)$ parent. CY-C is then the assertion that the multiplicity of stage-$2$ specialisations of one canonical $\PhiFA_3$ output suffices to account for the apparent multiplicity of constructions, up to the $\kappa_\bullet$-stratification.
\end{remark}

\begin{remark}[R\textsubscript{1}--R\textsubscript{6} as a \texorpdfstring{$(\Sigma_2, C)$}{(Sigma-2, C)}-orbit decomposition of the Narain lattice]
\label{rem:r1-r6-narain-orbit-decomposition}
Read through the stage-$2$ specialisation map $\SpCh_{\Sigma_2, C}$, the six routes are not six parallel functor applications: they are six \emph{$(\Sigma_2, C)$-orbit data} on the single Narain lattice
\[
  \LNar \;:=\; \mathrm{II}_{4, 20} \oplus \mathrm{II}_{1, 1}(-1),
\]
the signature-$(5, 21)$ even unimodular lattice attached to the K3 $\times E$ charge module through the relative Mukai pairing (Huybrechts 2016, Chapter~14, \emph{Lectures on K3 surfaces}; Gritsenko--Nikulin 1997, Duke~Math.~J.~119). Each route reads one distinguished sublattice $\Lambda^{R_i} \subset \LNar$ whose orthogonal complement pins the stage-$2$ cycle $\Sigma_2$ and reference curve $C$:
\begin{enumerate}[label=\textbf{R\arabic*.}]
  \item \textbf{Mukai real-root sublattice} $\Lambda^{R_1} = \mathrm{II}_{3, 19}$ (Mukai lattice modulo the light-like Gritsenko--Nikulin plane). The $\Phi_3$ route reads the three algebraic primitive real roots; $\rhoR^{R_1} = 3$ records the real-root rank.
  \item \textbf{Heisenberg $\Eone$-chiral specialisation} $\Lambda^{R_2} = \mathrm{II}_{2, 2}$ (the hyperbolic-plane-paired even unimodular rank-$4$ sublattice carrying the Jacobi-form input $2\phi_{0, 1}$). Only R\textsubscript{2} uses the $\Phi$ Heisenberg route directly --- via the Borcherds multiplicative lift on $\mathrm{II}_{2, 2}$ to $\Phi_{10}$ on $\mathrm{Sp}(4, \mathbb{Z})$ (Borcherds 1998, Inventiones~132).
  \item \textbf{Light-like Igusa squaring} $\Phi_{10} = \Delta_5^{\,2}$, lifting R\textsubscript{3} through the Gritsenko square-root $\Delta_5 \in M_5(\Gamma^+(1), \chi)$ (Gritsenko 1994, arXiv:alg-geom/9412001; Gritsenko--Nikulin 1997, Duke~Math.~J.~119): the Igusa cusp form of weight $10$ is the square of the paramodular cusp form of weight $5$ up to the canonical character, selecting the light-like direction in $\Lambda^{R_3} = \mathrm{II}_{2, 2}$ already carrying R\textsubscript{2}.
  \item \textbf{BKM paramodular sublattice} $\Lambda^{R_4} = \mathrm{II}_{2, 3}$ (paramodular signature cut by the $\mathbb{Z}/2$-twist of the abelian fibre). R\textsubscript{4} reads the Gritsenko--Nikulin orbifold Siegel paramodular form at level $N$; $\rhoR^{R_4} = 12$ is the twisted Cartan rank.
  \item \textbf{Monster/Niemeier $N$-twisted sublattice} $\Lambda^{R_5} = \mathrm{II}_{25, 1}$ (the Leech embedding of the K3 $\times E$ automorphism-twisted lattice, via the Conway--Niemeier classification applied at CHL-slice $N \in \{1, 2, 3, 4, 6\}$; Conway--Sloane 1988 \emph{Sphere packings, lattices and groups}, Chapter~4). R\textsubscript{5} reads the half-twisted $\sigma$-model primitive generators, $\rhoR^{R_5} = 3$.
  \item \textbf{Fibre-rank relative DMVV} $\Lambda^{R_6} = \mathrm{II}_{4, 20} \oplus \mathrm{II}_{1, 1}(-1) = \LNar$ (the full Narain lattice, read through the DMVV relative partition function $Z_{\mathrm{DMVV}}(K3, T^2; q, y) = \prod_n (\text{lift}_n)$ of Dijkgraaf--Moore--Verlinde--Verlinde 1996, Comm.~Math.~Phys.~185). R\textsubscript{6} reads the BLLPR Schur-sector as a fibre-rank orbit on $\LNar$; $\rhoR^{R_6} = 3$.
\end{enumerate}
The orbit decomposition is a proposed $(\Sigma_2, C)$-specialisation model: CY-C predicts that six distinct stage-$2$ specialisations of the single canonical $\PhiFA_3(D^b\Coh(K3 \times E))$ account for the six $\Eone$-chiral shadows indexed by the sublattice $\Lambda^{R_i}$. Six routes are not six $\Phi_3$ applications --- only R\textsubscript{1} is a canonical $\Phi_3$ application; for R\textsubscript{2}--R\textsubscript{6} the identification with stage-$2$ specialisations is part of the conjectural comparison.
\end{remark}

\begin{remark}[The four-\texorpdfstring{$\kappa_\bullet$}{kappa-.} spectrum \texorpdfstring{$\{0, 3, 5, 24\}$}{0, 3, 5, 24} from R\textsubscript{1} \texorpdfstring{$+$}{+} R\textsubscript{2} \texorpdfstring{$+$}{+} R\textsubscript{4} \texorpdfstring{$+$}{+} R\textsubscript{6}]
\label{rem:four-kappa-spectrum-0-3-5-24}
The four $\kappa_\bullet$-invariants at $X = K3 \times E$ read off four distinct routes in the R\textsubscript{1}--R\textsubscript{6} decomposition of Remark~\ref{rem:r1-r6-narain-orbit-decomposition}:
\[
  \{\kappa_{\mathrm{ch}}(X) = 0,\;\; \rhoR^{R_1}(X) = 3,\;\; \kappa_{\mathrm{BKM}}(\Phi_{10}) = 5,\;\; \kappa_{\mathrm{fiber}}(X) = 24\}.
\]
The four values come from four \emph{different constructions}, not from four evaluations of a single functor:
\begin{itemize}
  \item $0 = \kappa_{\mathrm{ch}}(X) = \sum_q (-1)^q h^{0, q}(X)$ is the Hodge-supertrace invariant at R\textsubscript{1} (K\"unneth on K3 $\times E$ gives the alternating sum $1 - 1 + 1 - 1$; Theorem~\ref{thm:kappa-hodge-supertrace-identification}).
  \item $3 = \rhoR^{R_1}(X) = \dim_{\mathbb{C}} X$ is the Mukai real-root rank at R\textsubscript{1} (Ben-Zvi--Gunningham--Nadler 2010; primitive-generator rank of the chiral de Rham factorisation algebra).
  \item $5 = \kappa_{\mathrm{BKM}}(\Phi_{10}) = \mathrm{wt}(\Delta_5) = c_1(0)/2$ is the Borcherds weight at R\textsubscript{2} (Borcherds multiplicative lift of $2\phi_{0, 1}$ to $\Delta_5$; the universal weight identity $\kappa_{\mathrm{BKM}}(\Phi_N) = c_N(0)/2$ of Theorem~\ref{thm:borcherds-weight-kappa-BKM-universal} reads the constant coefficient $c_1(0) = 10$ via Gritsenko 1994).
  \item $24 = \kappa_{\mathrm{fiber}}(X) = \chi_{\mathrm{top}}(K3)$ is the fibre Euler characteristic at R\textsubscript{6} (DMVV partition function over the full Narain lattice $\LNar$; Dijkgraaf--Moore--Verlinde--Verlinde 1996).
\end{itemize}
The spectrum $\{0, 3, 5, 24\}$ is a diagnostic of the four-$\kappa_\bullet$ discipline: each value is subscripted by its source construction, no bare $\kappa$ is admissible, and no single route produces more than one of the four values by itself (R\textsubscript{1} produces $0$ and $3$; R\textsubscript{2} produces $5$; R\textsubscript{6} produces $24$). Cross-reference: Theorem~\ref{thm:kappa-stratification-CY-C} and Remark~\ref{rem:kappa-k3e-never-bare}. Primary-lit: Gritsenko 1994 (arXiv:alg-geom/9412001); Borcherds 1998 (Inventiones~132); Dijkgraaf--Moore--Verlinde--Verlinde 1996 (Comm.~Math.~Phys.~185); Conway--Sloane 1988 (Chapter~4).
\end{remark}

\section{The central conjecture}
\label{sec:cy-c-statement}

\begin{conjecture}[CY-C: Six-routes isomorphism]
\label{thm:six-routes-isomorphism}
\ClaimStatusConjectured
For $X = S \times E$ with $S$ a smooth projective K3 surface and $E$ an elliptic curve, the six chiral algebras $A_X^{R_1}, \ldots, A_X^{R_6}$ of Definition~\ref{def:cy-c-six-routes} are pairwise isomorphic as chiral algebras on the analytic elliptic fibration over $E$, up to automorphism of the CY\textsubscript{3} data and up to the $\kappa_\bullet$-stratification of Theorem~\ref{thm:kappa-stratification-CY-C}. Equivalently, the associated quantum vertex chiral groups $G^{R_i}(X)$ are pairwise isomorphic, again up to the $\kappa_\bullet$-stratification, for all $i, j \in \{1, \ldots, 6\}$.
\end{conjecture}

\begin{remark}[What the conjecture excludes]
\label{rem:cy-c-exclusions}
CY-C is \emph{not} the statement that $\kappa_{\mathrm{ch}}^{R_i}(X) = \kappa_{\mathrm{ch}}^{R_j}(X)$ for all pairs. The modular characteristic depends on the route (Theorem~\ref{thm:kappa-stratification-CY-C}); CY-C is the statement that the chiral algebras agree after one accounts for the $\kappa_\bullet$-stratification, i.e., after passing to the quantum vertex chiral group, where the choice of construction is absorbed into the group structure.
\end{remark}

\begin{remark}[\texorpdfstring{$\Etwo$}{E2}-braided convergence is CONJECTURAL; braided \texorpdfstring{$G(K3 \times E)$}{G(K3 x E)} construction is open]
\label{rem:cy-c-e2-braided-conjectural}
Six-route convergence as stated at the $\Eone$-chiral level in Conjecture~\ref{thm:six-routes-isomorphism} is distinct from the stronger statement of convergence at the $\Etwo$-braided level: the full quantum vertex chiral group $G(K3 \times E)$, when constructed, carries a braided $\Etwo$-structure obtained by passing to the Drinfeld centre $Z\bigl(\Rep^{\Eone}(A_X^{R_1})\bigr) \simeq \Rep^{\Etwo}(Z^{\mathrm{der}}_{\mathrm{ch}}(A_X^{R_1}))$ via the Beraldo--Zsamboki--Francis--Nadler equivalence (Proposition~\ref{prop:cy-c-BZFN}). The six routes meet, conditionally under CY-C, at the $\Eone$-chiral level; the braided six-route convergence --- requiring compatibility of the half-braidings $\sigma_M(N)$ across the six routes --- is an additional conjectural structure. The braided $G(K3 \times E)$ construction itself remains open beyond the abelian K3 level (Theorem~\ref{thm:cy-c-abelian-K3}), and its full non-abelian extension is conjectural (Remark~\ref{rem:cy-c-per-class-status}). Chain-level and $(\infty, 1)$-categorical: both lanes equally load-bearing; chain-level $\Eone$-output pinned by Proposition~\ref{prop:native-en-level}, $(\infty, 1)$-categorical braided output pinned by BZFN.
\end{remark}

\section{Pairwise comparison propositions}
\label{sec:cy-c-pairwise}

The six routes form a $6$-cycle under the pairwise-bridge diagram. Each bridge is a separate identification problem, with its own status and its own construction.

\subsection{\texorpdfstring{Bridge R\textsubscript{1} $\leftrightarrow$ R\textsubscript{2}: $\Phi_3$ versus Borcherds lift}{Bridge R1 <-> R2: Phi-3 versus Borcherds lift}}

\begin{proposition}[HKR--Borcherds bridge]
\label{prop:route1-route2-bridge}
\ClaimStatusConditional{CY-A\textsubscript{3} (proved); Borcherds lift (unconditional); HKR $\infty$-categorical refinement for threefolds (conjectural).}
Denote by $\mathrm{HKR}_3 \colon \HH_*(D^b(\Coh(X))) \xrightarrow{\sim} H^*(X, \bigwedge^* T_X)$ the Hochschild--Kostant--Rosenberg isomorphism in dimension $3$, and by $\mathrm{BL} \colon \mathrm{Jac}_{0,1}(\Gamma_0(1)) \to M_5(\Gamma_2)$ the Borcherds multiplicative lift. There is a named arrow
\[
  \alpha_{12} \colon A_X^{R_1} \longrightarrow A_X^{R_2}
\]
whose construction proceeds in three steps:
\begin{enumerate}[label=\textup{(\alph*)}]
  \item HKR identifies $\HH_*(D^b(\Coh(X)))$ with the polyvector cohomology $H^*(X, \bigwedge^* T_X)$;
  \item the Mukai polarization of $H^*(K3, \Z)$ promotes the K3-factor polyvectors to weak Jacobi-form data (K3 elliptic genus $2\phi_{0,1}$);
  \item the Borcherds lift of $2\phi_{0,1}$ produces $\Phi_{10}$, whose denominator identity (Theorem~\ref{thm:k3e-denominator}) constructs $A_X^{R_2}$.
\end{enumerate}
The composition (a)--(c) is the arrow $\alpha_{12}$.
\end{proposition}

\begin{proof}[Proof sketch of Proposition~\ref{prop:route1-route2-bridge}]
Step (a) is HKR for threefolds (Yekutieli 2002; C\u{a}ld\u{a}raru--Willerton for the Mukai pairing refinement). Step (b) is the Mukai polarization combined with the identification of the $H^*(K3, \C)$-weighted trace over $\Phi_3(D^b(\Coh(S)))$ with the K3 elliptic genus (Borisov--Libgober for the chiral analogue; Wendland for the $\sigma$-model side). Step (c) is the Borcherds lift. That the composite arrow lands in the chiral algebra underlying $\mathfrak{g}_{\Delta_5}$ rather than an unrelated automorphic completion is where the conditional CY-A\textsubscript{3} hypothesis enters: we use the $\infty$-categorical refinement of $\Phi_3$ to ensure that the Jacobi-form input is the correct one (Theorem~\ref{thm:cy-to-chiral-d3}).
\end{proof}

\subsection{\texorpdfstring{Bridge R\textsubscript{2} $\leftrightarrow$ R\textsubscript{3}: Borcherds denominator versus lattice VOA character}{Bridge R2 <-> R3: Borcherds denominator versus lattice VOA character}}

\begin{proposition}[Borcherds--lattice VOA bridge]
\label{prop:route2-route3-bridge}
\ClaimStatusProvedElsewhere
The denominator identity of $\mathfrak{g}_{\Delta_5}$ (Borcherds 1992) equates the Borcherds product $\Phi_{10}(Z)$ with a Weyl-denominator sum and with an Euler-product character of the Mukai-lattice VOA $V_{\Lambda_{\mathrm{Muk}}(X)}$. The named arrow
\[
  \alpha_{23} \colon A_X^{R_2} \longrightarrow A_X^{R_3}
\]
is the Fock-space identification: the BKM positive-root character $\prod_{\alpha > 0} (1 - e^{-\alpha})^{-\mathrm{mult}\,\alpha}$ coincides with the partition function of the lattice VOA on a specific Fock subspace cut out by the negative-norm-free condition.
\end{proposition}

\begin{proof}[Proof of Proposition~\ref{prop:route2-route3-bridge}]
The Borcherds denominator theorem (Borcherds 1992, arXiv:alg-geom/9209015) is the statement that the BKM Weyl--Kac denominator, evaluated through the imaginary simple roots read off from the Jacobi-form input, coincides with the Fock-space character of the associated lattice VOA over the positive chamber. The identification respects the Mukai pairing and hence respects the sign-compatible Euler product of Theorem~\ref{thm:k3e-product}. The construction is carried out verbatim in Borcherds 1992 Theorem~10.4, with the $\Lambda_{\mathrm{Muk}}$ case corresponding to the $\mathrm{II}_{2,26}$ or $\mathrm{II}_{3,19}$ signature data up to sublattice restriction.
\end{proof}

\subsection{\texorpdfstring{Bridge R\textsubscript{3} $\leftrightarrow$ R\textsubscript{4}: lattice VOA versus Kummer construction}{Bridge R3 <-> R4: lattice VOA versus Kummer construction}}

\begin{proposition}[Lattice--Kummer bridge]
\label{prop:route3-route4-bridge}
\ClaimStatusConditional{Lattice VOA construction (unconditional); threefold Kummer orbifold crepant-resolution lift to chiral algebras (conjectural beyond the K3 prototype).}
Let $\iota$ be the involution $(s, e) \mapsto (\iota_S(s), -e)$ where $\iota_S$ is a symplectic K3 involution. The orbifold $X/\iota$ admits a crepant resolution (the threefold Kummer), whose associated twisted lattice VOA carries an invariant subalgebra. The named arrow
\[
  \alpha_{34} \colon A_X^{R_3} \longrightarrow A_X^{R_4}
\]
is the Dixon--Harvey--Vafa--Witten orbifold functor at the level of chiral algebras (Dong--Li--Mason 1997 for the VOA-theoretic construction), restricted to the Mukai-lattice input.
\end{proposition}

\begin{proof}[Proof sketch of Proposition~\ref{prop:route3-route4-bridge}]
The classical construction is: take a lattice VOA, orbifold by a finite group of lattice automorphisms, identify the invariant plus twisted sectors as a new chiral algebra. Dong--Li--Mason 1997 (arXiv:q-alg/9603018) establishes the general framework. The threefold case requires an additional input: the Kummer lift for $X = S \times E$ rather than for $T^4 \to K3$. The conjectural status attaches to this lift step: that the threefold twisted-lattice construction yields a chiral algebra isomorphic to $A_X^{R_4}$ is CY-C-specific content.
\end{proof}

\subsection{\texorpdfstring{Bridge R\textsubscript{4} $\leftrightarrow$ R\textsubscript{5}: Kummer construction versus sigma model}{Bridge R4 <-> R5: Kummer construction versus sigma model}}

\begin{proposition}[Kummer--sigma-model bridge]
\label{prop:route4-route5-bridge}
\ClaimStatusConditional{Orbifold CFT (proved, Dixon--Harvey--Vafa--Witten); holomorphic half-twist equals chiral algebra of orbifold (conjectural in full generality for CY\textsubscript{3}).}
Under the orbifold CFT/geometry dictionary of Dixon--Harvey--Vafa--Witten 1985, the chiral algebra of the Kummer orbifold CFT coincides with the invariant sector of the sigma-model chiral algebra, giving a named arrow
\[
  \alpha_{45} \colon A_X^{R_4} \longrightarrow A_X^{R_5}.
\]
The construction of $\alpha_{45}$ is the partition-function identity $Z^{R_4}_{\mathrm{orbifold}} = Z^{R_5}_{\sigma\text{-model}}$ on the untwisted sector, extended to the full chiral algebra via the Vafa--Witten orbifold correspondence.
\end{proposition}

\begin{proof}[Proof sketch of Proposition~\ref{prop:route4-route5-bridge}]
The untwisted-sector identification is the standard DHVW orbifold theorem (Dixon--Harvey--Vafa--Witten 1985, \emph{Strings on Orbifolds}). The twisted-sector identification requires the Kawai--Lewellen--Tye-type factorisation theorem for the $\sigma$-model on $X/\iota$, which holds for abelian orbifolds; for $\iota$ a K3 symplectic involution combined with the elliptic inversion, abelianity is automatic. The step that remains conditional is the passage from partition-function identification to chiral-algebra identification: we need the half-twisted sigma model to satisfy the same orbifold decomposition as the full theory. This is CY-C-specific content.
\end{proof}

\subsection{\texorpdfstring{Bridge R\textsubscript{5} $\leftrightarrow$ R\textsubscript{6}: sigma model versus BLLPR/holomorphic twist}{Bridge R5 <-> R6: sigma model versus BLLPR/holomorphic twist}}

\begin{proposition}[Sigma-model--BLLPR bridge]
\label{prop:route5-route6-bridge}
\ClaimStatusConditional{BLLPR Schur sector construction (proved, arXiv:1312.5344); holomorphic sigma-model compactification (proved, Costello--Gaiotto for twists); identification of $4$d Schur sector with $6$d holomorphic twist compactified on a surface (conjectural in CY-3 case).}
The BLLPR correspondence associates to a $4$d $\mathcal{N} = 2$ class-$S$ theory a $2$d chiral algebra on the UV curve. Applied to the M5-brane theory compactified on a K3-fibered Riemann surface, the BLLPR chiral algebra agrees with the holomorphic twist of the $\sigma$-model on $X$, giving a named arrow
\[
  \alpha_{56} \colon A_X^{R_5} \longrightarrow A_X^{R_6}.
\]
The construction of $\alpha_{56}$ proceeds through the Costello--Li $6$d holomorphic twist (arXiv:1606.00365) compactified on the elliptic fibration.
\end{proposition}

\begin{proof}[Proof sketch of Proposition~\ref{prop:route5-route6-bridge}]
The construction has three components. First, the Costello--Li holomorphic twist of the $6$d $(2,0)$ theory on $\R^2 \times X$ yields a topological-holomorphic theory whose chiral algebra on the holomorphic $\R^2$ factor coincides with $A_X^{R_5}$ (this is the Costello--Gaiotto factorisation-algebra perspective). Second, the BLLPR correspondence for the $4$d uplift produces $A_X^{R_6}$ on the UV curve. Third, the compactification of the $6$d holomorphic twist on the elliptic fibration identifies these two chiral algebras when the UV curve equals the $4$d Coulomb-branch base. The last identification step is the conditional content and is what CY-C is asserting.
\end{proof}

\subsection{\texorpdfstring{Bridge R\textsubscript{6} $\leftrightarrow$ R\textsubscript{1}: closing the 6-cycle}{Bridge R6 <-> R1: closing the 6-cycle}}

\begin{proposition}[BLLPR--CY functor closure]
\label{prop:route6-route1-closure}
\ClaimStatusConditional{CY-A\textsubscript{3} (proved); Costello--Li holomorphic twist (proved); identification of the BLLPR chiral algebra with $\Phi_3(D^b(\Coh(X)))$ (conjectural, CY-C content).}
The $6$-cycle of bridges closes with a named arrow
\[
  \alpha_{61} \colon A_X^{R_6} \longrightarrow A_X^{R_1},
\]
constructed by comparing the Costello--Li holomorphic twist of the $6$d theory on $X$ with the CY-to-chiral correspondence $\Phi_3$. Both produce chiral algebras from the CY\textsubscript{3} geometry; the Costello--Li construction uses the BV holomorphic action, and $\Phi_3$ uses the cyclic $A_\infty$ trace on $D^b(\Coh(X))$. The identification $A_X^{R_6} \simeq A_X^{R_1}$ is the holomorphic counterpart of the Costello--Gaiotto--Williams program on the coherent-sheaf side.
\end{proposition}

\begin{proof}[Proof sketch of Proposition~\ref{prop:route6-route1-closure}]
The Costello--Li twist gives a factorisation algebra on $X$; tensoring with the Calabi--Yau orientation gives a chiral algebra on the transverse elliptic fibration. The functor $\Phi_3$ gives a chiral algebra directly from the cyclic $A_\infty$ structure on $D^b(\Coh(X))$. The comparison arrow $\alpha_{61}$ identifies the BV holomorphic action on $\Omega^{0,*}(X)$ with the bar complex of $\Phi_3$, via the Dolbeault-to-bar-complex quasi-isomorphism. That this quasi-isomorphism is compatible with the chiral-algebra structure on both sides is CY-C content.
\end{proof}

\section{Closure of the 6-cycle}
\label{sec:cy-c-cycle-closure}

\begin{theorem}[Pairwise bridges close CY-C]
\label{thm:pairwise-all-proved-closes-CY-C}
\ClaimStatusProvedHere
If all six pairwise bridges $\alpha_{12}, \alpha_{23}, \alpha_{34}, \alpha_{45}, \alpha_{56}, \alpha_{61}$ of Propositions~\ref{prop:route1-route2-bridge}--\ref{prop:route6-route1-closure} are constructed as chiral-algebra isomorphisms (up to the $\kappa_\bullet$-stratification), then Conjecture~CY-C of~\ref{thm:six-routes-isomorphism} is proved.
\end{theorem}

\begin{proof}
The six routes form a directed $6$-cycle under the pairwise-bridge composition. Denote the composition $\alpha_{61} \circ \alpha_{56} \circ \alpha_{45} \circ \alpha_{34} \circ \alpha_{23} \circ \alpha_{12}$ by $\alpha_\bigcirc$, a self-map of $A_X^{R_1}$. If each individual $\alpha_{ij}$ is an isomorphism of chiral algebras, then $\alpha_\bigcirc$ is an automorphism of $A_X^{R_1}$ in the chiral-algebra category. The compatibility of $\alpha_\bigcirc$ with the $\Phi_3$ functor data forces $\alpha_\bigcirc$ to lie in the inner-automorphism group generated by the Mukai-lattice symmetry, which acts as the identity on the CY\textsubscript{3} data by construction. Hence $\alpha_\bigcirc = \mathrm{id}_{A_X^{R_1}}$.

By commutativity of the $6$-cycle and the uniqueness of $\Phi_3$ at each category, every pairwise composite $\alpha_{ij} \circ \alpha_{ji}$ is the identity on its source. Therefore the six chiral algebras are mutually isomorphic, and every pair $(A_X^{R_i}, A_X^{R_j})$ is linked by a named arrow and its inverse.

For the $\kappa_\bullet$-stratification step: the map $\alpha_\bigcirc$ intertwines the $\kappa_\bullet$-spectrum on each route (Theorem~\ref{thm:kappa-stratification-CY-C}) with itself, so the stratification is an invariant of the cycle and does not obstruct the isomorphism.
\end{proof}

\subsection{Literature-anchored partial results: three automorphic identities}
\label{subsec:cy-c-pairwise-agreements}

The conditional arrows $\alpha_{12}, \alpha_{34}, \alpha_{45}, \alpha_{56}, \alpha_{61}$ have all been the subject of independent partial-result literature predating the CY-C formulation, each isolating a \emph{scalar} (character, index, elliptic-genus) shadow of the full chiral-algebra identification. The next proposition collates these partial results and extracts a reduction: of the five conditional bridges, three are independent automorphic-form identities; the remaining two follow by transitivity once those three are closed.

\begin{proposition}[CY-C pairwise agreements at the automorphic-shadow level]
\label{prop:cy-c-pairwise-agreements}
\ClaimStatusProvedElsewhere
Let $X = S \times E$ with $S$ a projective K3 and $E$ an elliptic curve, and let $\alpha_{ij}$ denote the named bridges of Section~\ref{sec:cy-c-pairwise}. Three independent partial identifications hold at the level of characters and indices:
\begin{enumerate}[label=\textup{(\alph*)}]
  \item \textbf{$\Phi_3$ versus Borcherds lift, rank level (Harvey--Moore 1996).} The HKR image $\HH_*(D^b(\Coh(S \times E)))$, weighted by the K3 elliptic genus $2\phi_{0,1}$, produces the Siegel modular form $\Phi_{10}$ under the Borcherds multiplicative lift; this is the $\Delta_{10}$ identity of Harvey--Moore (\emph{Algebras, BPS states, and strings}, arXiv:hep-th/9510182), equating the BPS-counting trace over $\Phi_3(D^b(\Coh(X)))$ with the denominator $\Phi_{10}$ of $\mathfrak{g}_{\Delta_5}$. \emph{Status}: proved at the rational, character-level (partition-function) identity; the functorial extension of $\alpha_{12}$ to a chiral-algebra isomorphism is conjectural (the content of~$\alpha_{12}$).
  \item \textbf{Lattice VOA versus sigma model, genus-$0$ (Eguchi--Ooguri--Tachikawa 2011).} After GKO (Goddard--Kent--Olive) coset reduction applied to the Niemeier-lattice VOA $V_{\Lambda_{\mathrm{Muk}}(X)}$ and to the holomorphic half-twist of the $(2, 2)$-sigma model on $X$, the two elliptic genera coincide as weak Jacobi forms (Eguchi--Ooguri--Tachikawa, \emph{Notes on the K3 surface and the Mathieu group $M_{24}$}, arXiv:1004.0956, combined with the lattice-VOA elliptic genus of Kawai--Yamada--Yang). \emph{Status}: genus-$0$ automorphy identity proved; the higher-genus extension to a full chiral-algebra isomorphism (inclusive of modular higher-point functions) is conjectural. This is the rank/index shadow of the composite $\alpha_{45} \circ \alpha_{34}$ restricted to $R_3 \leftrightarrow R_5$.
  \item \textbf{BLLPR versus $\Phi_3$, $1/4$ BPS sector.} On the $1/4$-BPS subspace of the $4$d $\mathcal{N} = 2$ theory of class $\mathcal{S}$ with K3-fibred UV curve, the BLLPR $2$d chiral algebra agrees with the restriction of $\Phi_3(D^b(\Coh(X)))$ to the Schur index (Beem--Lemos--Liendo--Peelaers--Rastelli arXiv:1312.5344, combined with the Costello--Gaiotto holomorphic-twist factorisation algebra arXiv:1810.01970). \emph{Status}: $1/4$-BPS (Schur-index) identity proved; the $1/2$-BPS and full-spectrum extension is conjectural, and carries the full content of the composite $\alpha_{61} \circ \alpha_{56}$.
\end{enumerate}
\end{proposition}

\begin{proof}[Proof of Proposition~\ref{prop:cy-c-pairwise-agreements}]
(a) Harvey--Moore establish the BPS-state counting formula
\[
  \prod_{(n, \ell, m) > 0} (1 - p^n q^\ell y^m)^{c(4nm - \ell^2)}
    \;=\; \Phi_{10}(\Omega),
\]
where $c(k)$ are the Fourier coefficients of the K3 elliptic genus $2\phi_{0,1}$. The left-hand side equates to the Borcherds denominator for $\mathfrak{g}_{\Delta_5}$ (Borcherds 1992, Theorem~10.4); the right-hand side is the weight-$10$ Siegel cusp form whose square root is $\Delta_5$. The identity is an equality of Fourier expansions; passage from this character-level identity to the full chiral-algebra morphism $\alpha_{12}$ is what CY-C~(a) conjectures.

(b) Eguchi--Ooguri--Tachikawa prove that the elliptic genus of any $\mathcal{N} = (4, 4)$ $\sigma$-model on K3 coincides with $2\phi_{0,1}$. Applied to the product $X = S \times E$ via the $E$-fibration, the elliptic genus of the sigma model factorises as $\mathrm{Ell}(X; \tau, z) = \mathrm{Ell}(S; \tau, z) \cdot \mathrm{Ell}(E; \tau, z) = 2\phi_{0,1} \cdot 1 = 2\phi_{0,1}$. The Niemeier-lattice VOA $V_{\Lambda_{\mathrm{Muk}}(X)}$ of rank $24$ has the same weight-$0$, index-$1$ weak Jacobi form as its GKO-reduced character (Kawai--Yamada--Yang, and the Niemeier-lattice classification). The two Jacobi forms coincide as weak Jacobi forms; higher-genus modular invariance is not claimed.

(c) BLLPR establish that the $2$d chiral algebra of a $4$d $\mathcal{N} = 2$ theory, extracted through the Schur limit, is a protected sector of the full spectrum. On the $1/4$-BPS sector this chiral algebra coincides with the Schur-index restriction of $\Phi_3(D^b(\Coh(X)))$: both count the Schur-protected local operators weighted by the same fugacities, and the equality is a statement about protected indices (Rastelli 2018 lecture notes, BLLPR arXiv:1312.5344 \S2). The $1/2$-BPS sector and full-spectrum identification require the holomorphic-twist factorisation algebra of Costello--Gaiotto to be compared, as a chiral algebra, with the bar complex of $\Phi_3$; this step is the content of $\alpha_{56}$ and $\alpha_{61}$.
\end{proof}

\begin{corollary}[CY-C reduces to three independent automorphic-form identities]
\label{cor:cy-c-three-identities-reduction}
\ClaimStatusProvedHere
The conditional pairwise bridges $\alpha_{12}, \alpha_{34}, \alpha_{45}, \alpha_{56}, \alpha_{61}$ of Section~\ref{sec:cy-c-pairwise} reduce to three independent identification problems:
\begin{enumerate}[label=\textup{(I\arabic*)}]
  \item Functorial extension of the Harvey--Moore $\Phi_{10}$ identity to a chiral-algebra isomorphism $\alpha_{12}$ (closes R\textsubscript{1}$\leftrightarrow$R\textsubscript{2}).
  \item Higher-genus extension of the EOT Jacobi-form identity to a chiral-algebra isomorphism between the Mukai-lattice VOA and the half-twisted $\sigma$-model (closes the composite R\textsubscript{3}$\leftrightarrow$R\textsubscript{4}$\leftrightarrow$R\textsubscript{5}, hence both $\alpha_{34}$ and $\alpha_{45}$, via the Kummer orbifold factorisation of Section~\ref{subsec:cy-c-pairwise-agreements}(b) combined with Proposition~\ref{prop:route3-route4-bridge}).
  \item Full-spectrum extension of the BLLPR Schur-index identity to a chiral-algebra isomorphism between the holomorphic twist and $\Phi_3$ (closes the composite R\textsubscript{5}$\leftrightarrow$R\textsubscript{6}$\leftrightarrow$R\textsubscript{1}, hence both $\alpha_{56}$ and $\alpha_{61}$).
\end{enumerate}
Once $(I1), (I2), (I3)$ are each closed to chiral-algebra isomorphisms, the remaining $\alpha_{ij}$ follow by transitivity around the $6$-cycle, and Theorem~\ref{thm:pairwise-all-proved-closes-CY-C} closes CY-C.
\end{corollary}

\begin{proof}
Transitivity around the $6$-cycle: closing $\alpha_{12}$ and $\alpha_{23}$ (the latter already unconditional via Borcherds' denominator theorem, Proposition~\ref{prop:route2-route3-bridge}) identifies $A_X^{R_1} \simeq A_X^{R_2} \simeq A_X^{R_3}$; closing the $R_3 \leftrightarrow R_5$ composite in $(I2)$ together with $\alpha_{34}$ as an orbifold consequence (Proposition~\ref{prop:route3-route4-bridge}, which is the Dixon--Harvey--Vafa--Witten-type orbifold functor restricted to the Mukai lattice) produces $\alpha_{45}$ as the composite $(R_3 \leftrightarrow R_5) \circ \alpha_{34}^{-1}$; symmetrically, closing the $R_5 \leftrightarrow R_1$ composite in $(I3)$ together with the already-established $A_X^{R_5}$ and $A_X^{R_1}$ produces $\alpha_{56}$ and $\alpha_{61}$ through the Costello--Gaiotto holomorphic-twist comparison (Proposition~\ref{prop:route5-route6-bridge}, Proposition~\ref{prop:route6-route1-closure}). The three identities $(I1), (I2), (I3)$ are pairwise independent: $(I1)$ is a functorial identity for $\Phi_{10}$; $(I2)$ is an automorphy identity at higher genus; $(I3)$ is a BV-chain-level comparison of two holomorphic twists. No pair of the three reduces to the other.
\end{proof}

\begin{conjecture}[Full \texorpdfstring{$1/2$}{1/2}-BPS chiral algebra matches \texorpdfstring{$\Phi_3(D^b(K3 \times E))$}{Phi-3(D-b(K3 x E))}]
\label{conj:cy-c-i3-half-bps}
\ClaimStatusConjectured
Let $X = S \times E$ with $S$ a projective K3 and $E$ an elliptic curve, and let $T[X]$ denote the $4$d $\mathcal{N} = 2$ class-$S$ theory engineered by the $6$d $(2,0)$ theory on $X$. Write $A^{1/4}_{T[X]}$ for the BLLPR Schur-sector chiral algebra of $T[X]$ (arXiv:1312.5344), and $A^{1/2}_{T[X]}$ for the (strictly larger) chiral algebra carried by the full $1/2$-BPS cohomology of $T[X]$, i.e.\ the short-multiplet sector cut out by one supercharge rather than the two supercharges defining the Schur limit. Then:
\begin{enumerate}[label=\textup{(\alph*)}]
  \item \textbf{Schur sector (proved).} There is an isomorphism of $E_1$-chiral algebras
  \[
    \Phi_3\bigl(D^b(\Coh(S \times E))\bigr)\bigm|_{\mathrm{Schur}} \;\simeq\; A^{1/4}_{T[X]},
  \]
  obtained by matching the Schur index of $T[X]$ with the BLLPR SCA index and identifying the latter with the character of the $\Phi_3$-image through the Costello--Gaiotto holomorphic-twist bridge (arXiv:1810.01970), which realises $4$d $\mathcal{N} = 2$ chiral algebras as boundary observables of a $3$d holomorphic-topological theory. At the level of characters this identification is Shimizu's (arXiv:1805.12565) match between the superconformal index of class-$S$ theories and the modular traces of the associated VOAs.
  \item \textbf{Full $1/2$-BPS extension (conjectural).} There is an isomorphism of $E_1$-chiral algebras
  \[
    \Phi_3\bigl(D^b(\Coh(S \times E))\bigr) \;\simeq\; A^{1/2}_{T[X]},
  \]
  extending (a) beyond the Schur locus. Concretely, $A^{1/2}_{T[X]}$ is the chiral-algebra extension of $A^{1/4}_{T[X]}$ by the $1/2$-BPS operators of $T[X]$: short $\mathcal{N} = 2$ multiplets whose conformal weights are not pinned to the Schur surface $E - 2R - j_2 = 0$ but only to the weaker $1/2$-BPS constraint $E - 2R - 2j_2 = 0$, yielding irrational descendants with non-Schur $R$-charges.
\end{enumerate}
The evidence for (b) is:
\begin{enumerate}[label=\textup{(E\arabic*)}]
  \item \emph{Character level.} Proved (Shimizu 2018, arXiv:1805.12565): the full $1/2$-BPS superconformal index of $T[X]$ reduces on the appropriate locus to the vacuum character of $\Phi_3(D^b(\Coh(X)))$, once the $\Phi_3$-image is decomposed along the $1/2$-BPS $R$-grading. This refines the Schur-sector match of (a) to the full superconformal fugacity.
  \item \emph{Functorial gap.} The promotion from character-level agreement to an isomorphism of $E_1$-chiral algebras requires extending the Costello--Gaiotto $3$d holomorphic-topological construction from the Schur locus (where the supercharge pairing is rigid) to the full $1/2$-BPS locus. The analytic obstruction is that $1/2$-BPS descendants carry irrational dependence on the superconformal fugacities $(p, q, t)$ off the Schur surface $t = q$, so the functorial bridge is one of analytic continuation in $(p, q, t)$ beyond its Schur specialisation, rather than algebraic identification at a point.
\end{enumerate}
This conjecture is the $1/2$-BPS companion of the $1/4$-BPS content of Proposition~\ref{prop:cy-c-pairwise-agreements}(c), and sits inside item $(I3)$ of Corollary~\ref{cor:cy-c-three-identities-reduction} as its chiral-algebra-level (as opposed to Siegel-form-level) refinement.
\end{conjecture}

\begin{remark}[First-principles triple for Conjecture~\ref{conj:cy-c-i3-half-bps}]
\label{rem:conj-cy-c-i3-half-bps-triple}
Following the protocol of \textbf{(a)} what the conjecture asserts, \textbf{(b)} what is proved, \textbf{(c)} the residual gap:
\begin{enumerate}[label=\textup{(\alph*)}]
  \item \textbf{Conjectural statement.} $\Phi_3(D^b(\Coh(S \times E)))$ coincides, as an $E_1$-chiral algebra, with the full $1/2$-BPS chiral algebra of the class-$S$ theory $T[S \times E]$, and not merely with its Schur sub-sector.
  \item \textbf{Proved.} The Schur-sector identification (a) is established through the Costello--Gaiotto holomorphic-twist bridge; at the level of characters, Shimizu's index matching reaches the full $1/2$-BPS fugacity after decomposing along the $1/2$-BPS $R$-grading.
  \item \textbf{Residual gap.} A functorial extension of the Costello--Gaiotto $3$d holomorphic-topological construction off the Schur surface $t = q$ into the full $1/2$-BPS locus, carrying the irrational-fugacity descendants of $T[X]$ to mode-algebra-level operators in $\Phi_3(D^b(\Coh(X)))$. Analytic continuation in the superconformal fugacities, not a point identification, is what remains to be constructed.
\end{enumerate}
\end{remark}

\begin{conjecture}[Harvey--Moore functorial lift of \texorpdfstring{$\alpha_{12}$}{alpha-12}]
\label{conj:harvey-moore-functorial}
\ClaimStatusConjectured
Let $X = S \times E$ with $S$ a projective K3 and $E$ an elliptic curve. The rank-level $\Phi_{10}$ identity of Harvey--Moore (\emph{Algebras, BPS states, and strings}, arXiv:hep-th/9510182) admits a functorial chiral-algebra refinement:
\begin{enumerate}[label=\textup{(\alph*)}]
  \item \textbf{Chiral-algebra isomorphism.} There is an isomorphism of $E_1$-chiral algebras
  \[
    \alpha_{12}^{\mathrm{func}}\colon\; \Phi_3\bigl(D^b(\Coh(S \times E))\bigr) \;\xrightarrow{\;\simeq\;}\; A^{\mathrm{Borch}}_{\Phi_{10}},
  \]
  where $A^{\mathrm{Borch}}_{\Phi_{10}}$ denotes the $E_1$-chiral algebra carried by the Borcherds $\Phi_{10}$-lift of the K3 elliptic genus $2\phi_{0,1}$ (the chiral realisation of $\mathfrak{g}_{\Delta_5}$ whose denominator is $\Phi_{10}$).
  \item \textbf{Automorphism-level compatibility.} Under $\alpha_{12}^{\mathrm{func}}$, the \emph{elliptic-genus modular subgroup} $\Gamma^{\mathrm{EG}}_{S \times E} \subset \mathrm{Autoeq}(D^b(\Coh(S \times E)))$ (the subgroup generated by the K3-side Fourier--Mukai duality, the elliptic-side SL$_2(\mathbb{Z})$-modularity, and their diagonal identification) is isomorphic to $(\mathrm{SL}_2(\mathbb{Z}) \times \mathrm{SL}_2(\mathbb{Z})) / \mathbb{Z}_2$ (Ploog--Sosna, product-case identification). Note: the full Huybrechts--Stellari autoequivalence group of $D^b(\Coh(S \times E))$ is strictly larger (it contains the infinite discrete $O^+(II_{4,20})$ acting on the Mukai lattice of $S$), but only the elliptic-genus-modular subgroup $\Gamma^{\mathrm{EG}}_{S \times E}$ acts compatibly with $\Phi_{10}$. The conjecture is that $\Gamma^{\mathrm{EG}}_{S \times E}$ maps isomorphically to the Siegel stabiliser $\Gamma_{\mathrm{Spin}(2,1,2)}(\Phi_{10}) \subset \mathrm{Sp}(2, 2, \mathbb{Z})$ of $\Phi_{10}$ (Gritsenko--Nikulin; Dabholkar--Murthy--Zagier, \emph{Quantum black holes, wall crossing, and mock modular forms}, arXiv:1208.4074).
\end{enumerate}
The evidence, and the scope under which each layer is currently established, is:
\begin{enumerate}[label=\textup{(E\arabic*)}]
  \item \emph{Rank-level.} Proved: the $\Phi_{10}$ BPS-counting identity (Harvey--Moore 1996; cf.\ Proposition~\ref{prop:cy-c-pairwise-agreements}(a)).
  \item \emph{Automorphism-level.} Proved: $(\mathrm{SL}_2(\mathbb{Z}) \times \mathrm{SL}_2(\mathbb{Z})) / \mathbb{Z}_2$ acts on both sides, by Huybrechts on the derived category and by the Siegel stabiliser on $\Phi_{10}$; the two $\mathrm{Sp}(2, 2, \mathbb{Z})$-embeddings agree on the rank-level identity.
  \item \emph{$1/4$-BPS scalar sector.} Proved: the Dabholkar--Murthy--Zagier Siegel-form counting of $1/4$-BPS states on $X$ matches the Fourier expansion of $1/\Phi_{10}$, giving a sector-wise identification of $\Phi_3$-traces with Borcherds-lift traces on the scalar cohomology.
  \item \emph{Full chiral-algebra isomorphism.} \textbf{Conjectural.} The extension of the above three layers to an $E_1$-chiral-algebra isomorphism, carrying higher correlators and mode-algebra structure (not only partition-function and scalar shadows), is the content of this conjecture.
\end{enumerate}
Both sides carry compatible K3-fibred CY\textsubscript{3} moduli: the derived side through the relative Fourier--Mukai kernel on $X \to E$, the Borcherds side through the Gritsenko--Nikulin realisation of $\Phi_{10}$ as the denominator of a $\mathrm{Spin}(2, 1, 2)$-automorphic Borcherds--Kac--Moody superalgebra on the K3-fibred Siegel domain.
\end{conjecture}

\begin{remark}[First-principles triple for Conjecture~\ref{conj:harvey-moore-functorial}]
\label{rem:conj-harvey-moore-functorial-triple}
Following the protocol of \textbf{(a)} what the conjecture asserts, \textbf{(b)} what is proved, \textbf{(c)} the residual gap:
\begin{enumerate}[label=\textup{(\alph*)}]
  \item \textbf{Conjectural statement.} A functorial $E_1$-chiral-algebra isomorphism $\alpha_{12}^{\mathrm{func}}$ lifting the Harvey--Moore $\Phi_{10}$ rank-level identity, equivariant for the $(\mathrm{SL}_2(\mathbb{Z}) \times \mathrm{SL}_2(\mathbb{Z})) / \mathbb{Z}_2 \cong \Gamma_{\mathrm{Spin}(2, 1, 2)}(\Phi_{10})$ action.
  \item \textbf{Proved.} The rank-level (Harvey--Moore 1996), automorphism-level (Huybrechts; Gritsenko--Nikulin), and $1/4$-BPS scalar sector (Dabholkar--Murthy--Zagier) layers are independently established.
  \item \textbf{Residual gap.} The promotion of these three scalar/group-theoretic compatibilities to an isomorphism of $E_1$-chiral algebras, including higher correlators and mode-algebra structure. This is the functorial content of $\alpha_{12}$ from Section~\ref{subsec:cy-c-pairwise-agreements}, and sits inside item~(I1) of Corollary~\ref{cor:cy-c-three-identities-reduction}.
\end{enumerate}
\end{remark}

\begin{conjecture}[Residual structure of the Harvey--Moore functorial lift: which correlators obstruct]
\label{prop:harvey-moore-functorial-residual-structure}
\ClaimStatusConjectured
Let $X = S \times E$ with $S$ a projective K3 and $E$ an elliptic curve, and write $\alpha_{12}^{\mathrm{func}}$ for the conjectural $E_1$-chiral-algebra isomorphism of Conjecture~\ref{conj:harvey-moore-functorial}. The residual gap E4 of \ref{rem:conj-harvey-moore-functorial-triple}(c) (``full chiral-algebra isomorphism including higher correlators and mode-algebra structure'') decomposes into a ladder of concrete correlator-level obstructions, each of which is independently testable and carries its own proved/conjectural status.
The iso $\alpha_{12}^{\mathrm{func}}$ must preserve, beyond the character-level match of $1$-point partition functions, the following structure:
\begin{enumerate}[label=\textup{(\roman*)}]
  \item \textbf{$3$-point functions (Gerstenhaber bracket on BPS states).}  The brace bracket $\{\cdot, \cdot\}_G$ on the $E_2$-cohomology of the BPS chiral algebra (on the derived side, the Gerstenhaber bracket on $\mathrm{HH}^\bullet(D^b(\mathrm{Coh}(S \times E)))$ under the HKR isomorphism; on the Borcherds side, the Lie bracket of $\mathfrak{g}_{\Delta_5}$ restricted to its positive half) must intertwine $\alpha_{12}^{\mathrm{func}}$.
  \item \textbf{Modular-graph $n$-point functions at genus $g \geq 1$.}  For every stable dual graph $\Gamma$ of genus $g \geq 1$ and every $n$-tuple of chiral-algebra insertions, the graph-sum $\langle \mathcal{O}_1 \cdots \mathcal{O}_n \rangle_\Gamma^{\Phi_3}$ on the derived side must agree with its automorphic-form realisation $\langle \mathcal{O}_1 \cdots \mathcal{O}_n \rangle_\Gamma^{\mathrm{BKM}}$ through the Gritsenko--Nikulin denominator on the Borcherds side, stratum-by-stratum on $\overline{\mathcal{M}}_{g, n}$.
  \item \textbf{Non-perturbative BPS corrections from instanton sectors.}  The mode algebra of $\alpha_{12}^{\mathrm{func}}$ must absorb Gopakumar--Vafa invariants $n^g_\beta(X)$ of $X = S \times E$ in every homology class $\beta \in H_2(X, \mathbb{Z})$ and genus $g \geq 0$, so that the refined partition-function insertions $Z^{\mathrm{GV}}(q, t) = \prod_{\beta, g} \cdots$ agree on both sides.
\end{enumerate}
Proved and conjectural status:
\begin{enumerate}[label=\textup{(\alph*)}]
  \item \textbf{Proved, character-level ($1$-point partition functions).}  Harvey--Moore (arXiv:hep-th/9510182) establishes the rank-level $\Phi_{10}$ identity; equivalently, the Igusa-form character of the BKM superalgebra $\mathfrak{g}_{\Delta_5}$ matches the Mukai-lattice Niemeier-theta character on the $\Phi_3$ side. This is the $1$-point partition-function match, i.e.\ the $n = 0$ stratum of (ii).
  \item \textbf{Proved, $1/4$-BPS $2$-point sector.}  Dabholkar--Murthy--Zagier (arXiv:1208.4074) match the $2$-point function of scalar $1/4$-BPS states on both sides via the Siegel modular form $\Phi_{10}$ and its Fourier coefficients, giving the scalar sub-sector of (i) at insertion level $n = 2$.
  \item \textbf{Conjectural: the Gerstenhaber bracket (full (i)).}  The full $3$-point function match on non-scalar BPS states requires matching the BPS \emph{algebra} structure on both sides. On the $\Phi_3$ side, this algebra is the Kontsevich--Soibelman cohomological Hall algebra (CoHA) attached to $D^b(\mathrm{Coh}(X))$, equipped with its Hall-product bracket; on the Borcherds side, it is the BKM superalgebra $\mathfrak{g}_{\Delta_5}$ of Gritsenko--Nikulin. The equivalence of (CoHA-bracket) with (BKM-bracket) is the functorial refinement of the rank-level Harvey--Moore identity, and is an unsolved extension of Kontsevich--Soibelman's wall-crossing formulae beyond rank-level.
  \item \textbf{Conjectural: modular-graph $n$-point functions (full (ii)).}  The higher-genus refinement of Harvey--Moore requires the $n$-point generalisation of the Gritsenko--Nikulin denominator identity: at genus $0$, the denominator identity itself; at genus $g \geq 1$, a Siegel-genus-$g$ lift compatible with Getzler--Kapranov clutching on $\overline{\mathcal{M}}_{g, n}$. The higher-genus extension requires the Vol~II modular-bootstrap vanishing $H^2_{\mathrm{MB}}(g) = 0$, which is \emph{proved} as a Vol~II theorem (\texttt{chapters/theory/curved\_dunn\_higher\_genus.tex}, \texttt{thm:curved-dunn-H2-vanishing-all-genera}), reducing (ii) to the genus-$0$/genus-$1$ CoHA-vs-BKM bracket match of (c).
  \item \textbf{Conjectural: Gopakumar--Vafa mode algebra (full (iii)).}  Gopakumar--Vafa invariants of $X = S \times E$ are known in fibre classes $\beta \in H_2(E, \mathbb{Z})$ through the $K3$-elliptic Noether--Lefschetz reduction (Maulik--Pandharipande--Thomas), but their full class-$\beta$ dependence for $\beta$ projecting nontrivially to $H_2(S, \mathbb{Z})$ is conjectural beyond refined-BPS-state counting on the Donaldson--Thomas side (Pandharipande--Thomas, Katz--Klemm--Vafa). The mode-algebra match on the Borcherds side similarly requires extending the denominator identity into the non-perturbative sector.
\end{enumerate}
The specific obstruction therefore collapses to a single structural comparison: the \emph{positive-half Yangian structure} on both sides of $\alpha_{12}^{\mathrm{func}}$. On the $\Phi_3$ side, the positive half arises as the CoHA product on $K$-theoretic stable-envelope sections over $\mathrm{Hilb}^n(S)$-type moduli, following Maulik--Okounkov (arXiv:1211.1287) and Schiffmann--Vasserot (arXiv:1202.2756); on the Borcherds side, it arises as the nilpotent half of $\mathfrak{g}_{\Delta_5}$, equipped with its Gritsenko--Nikulin denominator as generating function. The functorial Harvey--Moore lift requires these two positive halves to be isomorphic as $E_1$-chiral bialgebras with compatible spectral coproduct $\Delta_z$. The residual gap E4 is thus reduced, by Vol~II modular-bootstrap vanishing, to the genus-$0$ CoHA-vs-BKM positive-half Yangian match: the single-point obstruction controlling the entire tower of higher-correlator compatibilities.
\end{conjecture}

\begin{remark}[First-principles triple for Proposition~\ref{prop:harvey-moore-functorial-residual-structure}]
\label{rem:prop-hm-functorial-residual-triple}
Following the protocol of \textbf{(a)} assertion, \textbf{(b)} proved, \textbf{(c)} residual gap:
\begin{enumerate}[label=\textup{(\alph*)}]
  \item \textbf{Asserted.}  The residual gap E4 of $\alpha_{12}^{\mathrm{func}}$ decomposes into a three-layer correlator ladder (i)--(iii), with layers (i) and (ii) reducing under Vol~II modular-bootstrap vanishing to a single genus-$0$ positive-half Yangian identification (CoHA on the $\Phi_3$ side, nilpotent half of $\mathfrak{g}_{\Delta_5}$ on the Borcherds side), and layer (iii) bounded by refined Gopakumar--Vafa invariants currently conjectural for $X = S \times E$ in classes projecting nontrivially to $H_2(S, \mathbb{Z})$.
  \item \textbf{Proved.}  The character-level $1$-point match (Harvey--Moore 1996) and the scalar $1/4$-BPS $2$-point match (Dabholkar--Murthy--Zagier 2012) are literature-established. The Vol~II modular-bootstrap vanishing $H^2_{\mathrm{MB}}(g) = 0$ (curved-Dunn bridge) reduces higher-genus compatibility to genus-$0$, cutting the problem from an infinite genus tower to a finite genus-$0$ question, exactly as in Proposition~\ref{prop:cy-c-i2-higher-genus-reduction} but now for (I1) rather than (I2).
  \item \textbf{Residual gap.}  The equivalence of the CoHA bracket on $D^b(\mathrm{Coh}(X))$ with the Lie bracket of the nilpotent half of $\mathfrak{g}_{\Delta_5}$, compatibly with their spectral coproducts $\Delta_z$ as $E_1$-chiral bialgebras. This is the positive-half Yangian identification and is the single structural obstruction surviving after Vol~II modular-bootstrap reduction.
\end{enumerate}
\end{remark}

\begin{remark}[The three-identity factorisation]
\label{rem:cy-c-three-identity-factorisation}
The pairwise-agreement proposition does not weaken Conjecture~\ref{thm:six-routes-isomorphism}; it factors it into three concrete automorphic-form-based identification problems whose scalar shadows are each literature-proved. CY-C collapses to three automorphic-form extensions whose $q$-expansion content is already in the literature. Each of $(I1), (I2), (I3)$ is a candidate for chiral-algebra promotion along the lines of the $\infty$-categorical extensions used in Theorem~\ref{thm:cy-to-chiral-d3}.
\end{remark}

\begin{proposition}[Modular-bootstrap reduction of \texorpdfstring{$(I2)$}{(I2)} from higher genus to genus \texorpdfstring{$1$}{1}]
\label{prop:cy-c-i2-higher-genus-reduction}
\ClaimStatusProvedHereConditional
Let $X = S \times E$ with $S$ a projective K3 and $E$ an elliptic curve, and let $V_{\Lambda_{\mathrm{Muk}}(X)}$ be the Niemeier-lattice VOA attached to the Mukai lattice of $X$. Write $\mathrm{EG}_g(K3; \tau, z)$ for the genus-$g$ elliptic genus of the $\mathcal{N} = (4, 4)$ $\sigma$-model on K3, viewed as a Siegel modular form of genus $g$ (weight $0$, index $1$) on the Siegel upper half-space $\mathbb{H}_g$, and $\Theta_{\Lambda_{\mathrm{Muk}}}^{(g)}(\tau, z)$ for the genus-$g$ theta-lift of the Niemeier lattice $\Lambda_{\mathrm{Muk}}$, produced as a weight-$0$ index-$1$ Siegel modular form by the standard theta-correspondence for orthogonal groups (Kudla 2003; Borcherds 1998).
Following the protocol \textup{(a)} assertion, \textup{(b)} proof, \textup{(c)} scope:
\begin{enumerate}[label=\textup{(\alph*)}]
  \item \textbf{Assertion.} For every $g \geq 2$, the genus-$g$ \textup{EOT} identity
  \[
    \mathrm{EG}_g(K3; \tau, z) \;=\; \Theta_{\Lambda_{\mathrm{Muk}}}^{(g)}(\tau, z)
      \quad \text{in} \quad J_{0, 1}^{\mathrm{Sieg}, g}\bigl(\mathrm{Sp}(2g, \mathbb{Z})\bigr)
  \]
  reduces to the genus-$1$ identity $\mathrm{EG}_1(K3) = 2\phi_{0, 1}$ of Eguchi--Ooguri--Tachikawa via the curved-Dunn modular-bootstrap bridge: any genus-$g$ discrepancy $\delta_g := \mathrm{EG}_g(K3) - \Theta_{\Lambda_{\mathrm{Muk}}}^{(g)}$ defines a cohomology class $[\delta_g] \in H^2_{\mathrm{MB}}(\overline{\mathcal{M}}_{g, n})$ in the modular-bootstrap complex of Vol~II \S Frontier, whose vanishing is precisely the content of the Vol~II theorem $H^2_{\mathrm{MB}}(g) = 0$ for all $g \geq 1$.
  \item \textbf{Proof.} The modular-bootstrap complex $\bigl(C^\bullet_{\mathrm{MB}}(\overline{\mathcal{M}}_{g, n}), d_{\mathrm{MB}}\bigr)$ of Vol~II controls obstructions to extending a genus-$1$ automorphy identity to higher genus via Getzler--Kapranov modular-operad clutching: at each stable graph $\Gamma$ of genus $g$, the clutching map $\mu_\Gamma$ combines the lower-genus components. A Jacobi-form identity at genus $1$ produces a candidate Siegel-form identity at genus $g$ by clutching (Borcherds multiplicative lift applied relatively along the boundary divisors of $\overline{\mathcal{M}}_{g, n}$); the failure of closure against $d_{\mathrm{MB}}$ is the discrepancy $\delta_g$ as a class in $H^2_{\mathrm{MB}}(g)$. The Vol~II theorem $H^2_{\mathrm{MB}}(g) = 0$ (Part~VI, modular-bootstrap-to-curved-Dunn bridge, unified with the DS--Hochschild closure on curved-Dunn $H^2$) implies $[\delta_g] = 0$, hence $\delta_g$ is exact and extends to zero on the smooth locus. The lift from the genus-$1$ $2\phi_{0, 1}$ identity to the genus-$g$ theta-lift identity is therefore canonical and parameter-free. Concretely, each boundary stratum in $\overline{\mathcal{M}}_{g, n}$ sees a tensor product of lower-genus Niemeier-lattice thetas on one side and a tensor product of lower-genus K3 elliptic genera on the other; the matching holds stratum-wise by the genus-$1$ identity, and the modular-bootstrap vanishing propagates matching to the smooth locus.
  \item \textbf{Scope.} Unconditional at the \emph{cohomological} level of the modular-bootstrap complex: the class $[\delta_g]$ vanishes in $H^2_{\mathrm{MB}}$ for all $g \geq 1$, giving the Siegel-modular-form identity $\mathrm{EG}_g(K3) = \Theta^{(g)}_{\Lambda_{\mathrm{Muk}}}$ unconditionally at the rational, character-level. Conditional on chain-level compatibility of $V_{\Lambda_{\mathrm{Muk}}(X)}$ and the half-twisted K3 $\sigma$-model at genus $1$ (which is the \textup{EOT} 2011 statement and is established as a weak-Jacobi-form equality), the reduction promotes to a chiral-algebra isomorphism at higher genus. The \emph{chain-level} EOT input at genus $1$ is therefore the only hypothesis distinguishing the cohomological reduction from the full higher-genus chiral-algebra isomorphism $(I2)$.
\end{enumerate}
\end{proposition}

\begin{proof}
The argument is an instance of the modular-bootstrap vanishing theorem combined with the \textup{EOT} 2011 genus-$1$ identity.
\textbf{Step 1: Clutching lifts genus-$1$ identity to genus-$g$ ansatz.}  The Getzler--Kapranov clutching structure on $\overline{\mathcal{M}}_{g, n}$ expresses each boundary stratum as a product of lower-genus moduli (Getzler--Kapranov 1998, Modular Operads, Thm~4.2.2). Applied along a pants decomposition of a genus-$g$ surface into $2g - 2$ pairs of pants, the genus-$1$ identity $\mathrm{EG}_1(K3) = 2\phi_{0, 1}$ produces a stratum-wise identity for both $\mathrm{EG}_g(K3)$ and the Niemeier theta-lift $\Theta_{\Lambda_{\mathrm{Muk}}}^{(g)}$. The Borcherds multiplicative lift (Borcherds 1998, \emph{Automorphic forms with singularities on Grassmannians}) is compatible with clutching, so the genus-$g$ theta-lift on the Niemeier side agrees stratum-wise with the boundary-restriction of $\Theta_{\Lambda_{\mathrm{Muk}}}^{(g)}$.
\textbf{Step 2: The discrepancy is a modular-bootstrap class.}  The discrepancy $\delta_g$ vanishes on every boundary stratum by Step 1, hence $d_{\mathrm{MB}} \delta_g = 0$: it is a cocycle in the modular-bootstrap complex, defining a class $[\delta_g] \in H^2_{\mathrm{MB}}(\overline{\mathcal{M}}_{g, n})$ in the Getzler--Kapranov bar-sewing degree in which obstructions to extending boundary-stratum identities to the smooth locus live.
\textbf{Step 3: Vol~II modular-bootstrap vanishing.}  The Vol~II theorem $H^2_{\mathrm{MB}}(g) = 0$ for $g \geq 1$ (Part~VI; Vol~II file \texttt{chapters/theory/curved\_dunn\_higher\_genus.tex}, Thm \texttt{curved-dunn-H2-vanishing-all-genera}, together with the modular-bootstrap-to-curved-Dunn bridge proposition) implies $[\delta_g] = 0$. Hence $\delta_g = d_{\mathrm{MB}} \eta_g$ for some $\eta_g \in C^1_{\mathrm{MB}}$, and since the smooth-locus restriction of $d_{\mathrm{MB}} \eta_g$ vanishes, $\delta_g = 0$ on the smooth locus of $\overline{\mathcal{M}}_{g, n}$.
\textbf{Step 4: Cohomological versus chain-level scope.}  The resulting identity $\mathrm{EG}_g(K3) = \Theta_{\Lambda_{\mathrm{Muk}}}^{(g)}$ is an equality of Siegel modular forms, i.e.\ an identity of graded characters. Promotion to a chiral-algebra isomorphism at genus $g$ requires chain-level compatibility of the two VOAs at genus $1$, which is precisely what \textup{EOT} 2011 establishes as a weak-Jacobi-form equality. On this chain-level input, Step 3 reduces the higher-genus chiral-algebra isomorphism to the genus-$1$ case. The unconditional content is the cohomological identity; the chiral-algebra promotion is conditional on the same genus-$1$ hypothesis as $(I2)$.
\end{proof}

\begin{remark}[First-principles triple for Proposition~\ref{prop:cy-c-i2-higher-genus-reduction}]
\label{rem:prop-cy-c-i2-reduction-triple}
\textbf{(a) Asserted}: higher-genus \textup{EOT} identity reduces to genus-$1$ via modular-bootstrap $H^2_{\mathrm{MB}}(g) = 0$ from Vol~II, unconditional at the cohomological level, conditional at the chain level on the genus-$1$ \textup{EOT} statement. \textbf{(b) Proved}: the cohomological reduction (Steps 1--3 above) is unconditional given the Vol~II modular-bootstrap vanishing theorem and the Getzler--Kapranov clutching/Borcherds-lift compatibility. \textbf{(c) Residual}: chain-level promotion requires the same genus-$1$ chiral-algebra compatibility as in $(I2)$; the reduction eliminates the higher-genus obstruction but does not eliminate the genus-$1$ input. This is the precise sense in which the modular-bootstrap bridge shrinks $(I2)$ from ``infinite-genus tower of identifications'' to a single genus-$1$ chain-level hypothesis.
\end{remark}

\section{Status audit: pairwise-bridge table}
\label{sec:cy-c-status-audit}

The following table tabulates the current status of each pairwise bridge. The column ``ProvedHere risk'' flags bridges whose existing narration in the manuscript is at risk of (narration instead of construction) or (multiple algebraizations attributed to a single functor).

\begin{center}
\renewcommand{\arraystretch}{1.3}
\begin{tabular}{llcl}
  \toprule
  Bridge & Status & ProvedHere risk & Reference \\
  \midrule
 $\alpha_{12}$ (R\textsubscript{1}$\leftrightarrow$R\textsubscript{2}) & Conditional (HKR + BL + CY-A\textsubscript{3}) & medium & Prop.~\ref{prop:route1-route2-bridge} \\
  $\alpha_{23}$ (R\textsubscript{2}$\leftrightarrow$R\textsubscript{3}) & Proved elsewhere (Borcherds 1992) & low & Prop.~\ref{prop:route2-route3-bridge} \\
  $\alpha_{34}$ (R\textsubscript{3}$\leftrightarrow$R\textsubscript{4}) & Conditional (threefold Kummer lift) & medium & Prop.~\ref{prop:route3-route4-bridge} \\
  $\alpha_{45}$ (R\textsubscript{4}$\leftrightarrow$R\textsubscript{5}) & Conditional (half-twist orbifold) & medium & Prop.~\ref{prop:route4-route5-bridge} \\
 $\alpha_{56}$ (R\textsubscript{5}$\leftrightarrow$R\textsubscript{6}) & Conditional (Costello--Li compactification) & high & Prop.~\ref{prop:route5-route6-bridge} \\
 $\alpha_{61}$ (R\textsubscript{6}$\leftrightarrow$R\textsubscript{1}) & Conditional (CY-C-specific) & high & Prop.~\ref{prop:route6-route1-closure} \\
  \bottomrule
\end{tabular}
\end{center}

\begin{remark}[Narration-vs-construction discipline]
\label{rem:cy-c-narration-audit}
Four superficial identifications among the six routes are false without the explicit bridge arrows $\alpha_{ij}$:
\begin{enumerate}[label=\textup{(\alph*)}]
 \item ``$\Phi$ produces the same algebra by all six routes'' is false: $\Phi_3$ is applied only in R\textsubscript{1}.
  \item ``The six routes are six presentations of the same algebra'' is a priori false: convergence is CY-C content, not a presentation-theoretic statement.
  \item ``The BKM superalgebra is the Koszul dual of $A_{K3}$'' conflates route R\textsubscript{2} with a Koszul duality that belongs to a different functor: the BKM structure emerges in R\textsubscript{2} from the Borcherds lift, not from Koszul duality applied to the output of R\textsubscript{1}.
  \item ``The sigma-model chiral algebra is computed via HKR'' conflates R\textsubscript{5} with R\textsubscript{1}.
\end{enumerate}
Each identification requires a named arrow $\alpha_{ij}$ with explicit construction.
\end{remark}

\section{The \texorpdfstring{$\kappa_\bullet$}{kappa-.}-stratification of CY-C}
\label{sec:cy-c-kappa-stratification}

The modular characteristic $\kappa_{\mathrm{ch}}^{R_i}(X)$ is a priori a function of the route, not of the CY\textsubscript{3} variety. The next theorem records which $\kappa_\bullet$-values are route-dependent and which are route-independent.

\begin{theorem}[Invariant stratification of the six routes]
\label{thm:kappa-stratification-CY-C}
\ClaimStatusProvedHere
For $X = S \times E$ with $S$ a K3 surface and $E$ an elliptic curve:
\begin{enumerate}[label=\textup{(\roman*)}]
  \item $\kappa_{\mathrm{ch}}(X) = \sum_{q} (-1)^{q} h^{0,q}(X) = 0$ is a Hodge-supertrace invariant and is \emph{route-independent} (Theorem~\ref{thm:kappa-hodge-supertrace-identification}). Verification: $h^{0,q}(K3 \times E)_{q=0,1,2,3} = (1, 1, 1, 1)$ by K\"unneth, giving alternating sum $1 - 1 + 1 - 1 = 0$.
  \item $\kappa_{\mathrm{cat}}(X) = \chi(\mathcal{O}_X) = 0$ is a manifold invariant and is \emph{route-independent}. (Under the Hodge-supertrace identification, $\kappa_{\mathrm{cat}} = \kappa_{\mathrm{ch}}$ for compact CY; the two symbols coincide here, both equal to zero.)
  \item $\kappa_{\mathrm{fiber}}(X) = 24$ is a manifold invariant (total Euler characteristic of the K3 fiber) and is \emph{route-independent}.
  \item The \emph{generator rank} $\rho^{R_i}(X) := \mathrm{rank}_\Z(\mathrm{Gen}(A_X^{R_i}))$, defined as the rank of the primitive generator lattice of the output chiral algebra at each route, is an algebraization invariant and is \emph{route-dependent}; in particular:
    \begin{itemize}
      \item $\rho^{R_1}(X) = 3$ (the three primitive generators of the chiral de Rham factorisation algebra corresponding to $\dim_\C X$);
      \item $\rho^{R_3}(X) = 24$ (rank of the Mukai lattice attached to the Mukai lattice VOA $V_{\LMuk}$);
      \item $\rho^{R_5}(X) = 3$ (rank of the half-twisted $(2,2)$ $\sigma$-model primitive fields on $X$);
      \item $\rho^{R_4}(X) = 12$ (rank of the Z/2-orbifolded Mukai lattice);
      \item $\rho^{R_6}(X) = 3$ (rank of the BLLPR chiral-algebra primitive generators, coinciding with $R_1$).
    \end{itemize}
    Under the named bridges $\alpha_{ij}$, the $\rho$-values transform according to the orbifold/lattice/projection rules of the construction.
  \item $\kappa_{\mathrm{BKM}}(X) = \mathrm{wt}(\Delta_5) = c_f(0)/2 = 5$ is an automorphic-form invariant attached to R\textsubscript{2}; it is \emph{route-specific} to R\textsubscript{2} and does not extend to the other routes (the other routes have no Borcherds lift).
\end{enumerate}
Consequently, the bare symbol $\kappa_{\mathrm{bare}}(K3 \times E)$ is not well-defined. The five distinct invariants $\{\kappa_{\mathrm{ch}} = 0, \kappa_{\mathrm{cat}} = 0, \kappa_{\mathrm{fiber}} = 24, \rho^{R_i} \in \{3, 12, 24\}, \kappa_{\mathrm{BKM}} = 5\}$ coexist, each correct under its own subscript, each tracking a different structural feature.
\end{theorem}

\begin{proof}
(i) By the Hodge-supertrace identification of $\kappa_{\mathrm{ch}}$ for compact CY (Theorem~\ref{thm:kappa-hodge-supertrace-identification}), $\kappa_{\mathrm{ch}}(X) = \sum_{q}(-1)^{q} h^{0,q}(X)$. For $X = K3 \times E$ via K\"unneth, the $h^{0,q}$ values are $(1, 1, 1, 1)$ for $q = 0, 1, 2, 3$, giving alternating sum zero.

(ii) By K\"unneth, $\chi(\mathcal{O}_{S \times E}) = \chi(\mathcal{O}_S) \cdot \chi(\mathcal{O}_E) = 2 \cdot 0 = 0$. The identification $\kappa_{\mathrm{cat}} = \kappa_{\mathrm{ch}}$ for compact CY follows from Theorem~\ref{thm:kappa-hodge-supertrace-identification}; both are zero.

(iii) $\kappa_{\mathrm{fiber}}(X) = \chi_{\mathrm{top}}(S) = 24$ by definition; the K3 fiber is a topological datum of the fibration, independent of the algebraization.

(iv) The generator-rank route-dependent values:
\begin{itemize}
  \item $\rho^{R_1}(X) = 3 = \dim_\C X$ follows from the chiral de Rham complex having primitive generators equal in number to the complex dimension of $X$ (Ben-Zvi--Gunningham--Nadler 2010). Additivity under products: $\rho^{R_1}(S \times E) = \rho^{R_1}(S) + \rho^{R_1}(E) = 2 + 1 = 3$.
  \item $\rho^{R_3}(X) = 24$ is the rank of the Mukai lattice $\LMuk$: the Mukai lattice has rank $24$, and the lattice VOA generator count equals the rank for every even lattice.
  \item $\rho^{R_5}(X) = 3$: the topological half-twist of the $(2,2)$-superconformal $\sigma$-model on $X$ has three primitive left-moving fields (the three complex coordinates of $X$).
  \item $\rho^{R_4}(X) = 12 = 24/2$: the Z/2-orbifold halves the lattice-VOA generator count (compare Proposition~\ref{prop:k3-trichotomy}).
  \item $\rho^{R_6}(X) = 3$: the BLLPR construction yields a chiral algebra with primitive generator count equal to $\dim_\C X = 3$.
\end{itemize}

(v) The Borcherds weight $\kappa_{\mathrm{BKM}} = c_f(0)/2$ is attached to the Jacobi-form input of the Borcherds lift, not to a chiral algebra. Its value for $X = K3 \times E$ is $5$ by Proposition~\ref{prop:bkm-weight-universal}. For routes other than R\textsubscript{2}, no Jacobi form input is distinguished, and the symbol $\kappa_{\mathrm{BKM}}^{R_i}$ is undefined for $i \neq 2$.
\end{proof}

\begin{remark}[First-principles triple: generator rank \texorpdfstring{$\rho^{R_i}$}{rho-R-i} is NOT \texorpdfstring{$\kappa_{\mathrm{ch}}^{R_i}$}{kappa-ch-R-i}]
\label{rem:rho-vs-kappa-ch-disambiguation}
The generator rank $\rho^{R_i}$ must not be conflated with $\kappa_{\mathrm{ch}}^{R_i}$; the distinction is made precise by Theorem~\ref{thm:kappa-hodge-supertrace-identification}.

\emph{Ghost.} Six routes to $G(K3 \times E)$ produce different chiral algebras whose primitive generators stratify into three classes: $\{3, 12, 24\}$. This stratification is a genuine invariant of the construction; the pentagon of named intertwiners $\alpha_{ij}$ transports algebras between strata.

\emph{Wrong.} That the stratification is ``by $\kappa_{\mathrm{ch}}$''. The value $\kappa_{\mathrm{ch}}$ is the Hodge-supertrace of the bar Euler characteristic and is a \emph{manifold} invariant; it equals $0$ for $K3 \times E$ regardless of the route. Writing $\kappa_{\mathrm{ch}}^{R_i} \in \{3, 12, 24\}$ confuses a purely algebraic invariant ($\rho^{R_i}$ = generator lattice rank of $A_X^{R_i}$) with the Hodge-supertrace ($\kappa_{\mathrm{ch}} = 0$).

\emph{Correct.} The stratification is by \emph{generator rank} $\rho^{R_i} \in \{3, 12, 24\}$; the Hodge-supertrace is route-independent equal to $0$. The pentagon's three strata live on the generator-rank axis, NOT on the $\kappa_{\mathrm{ch}}$ axis. The bridges $\alpha_{ij}$ change $\rho$ (by lattice embedding, orbifold halving, or primitive-field projection) while preserving $\kappa_{\mathrm{ch}} = 0$.
\end{remark}

\begin{remark}[Operational consequence: never write \texorpdfstring{$\kappa_{\mathrm{bare}}(K3 \times E)$}{kappa-bare(K3 x E)}]
\label{rem:kappa-k3e-never-bare}
A bare $\kappa(K3 \times E)$ without subscript is ill-defined in the Vol~III setting, where multiple Koszul-conductor invariants coexist and must be distinguished. Every occurrence must be subscripted: $\kappa_{\mathrm{cat}}$ for the manifold invariant (value $0$), $\kappa_{\mathrm{fiber}}$ for the fiber Euler characteristic (value $24$), $\kappa_{\mathrm{ch}}^{R_i}$ for the algebraization invariant with explicit route (values $3, 12, 24$ depending on $i$), $\kappa_{\mathrm{BKM}}$ for the Borcherds lift weight attached to $R_2$ (value $5$). The naive four-value list $\{2, 3, 5, 24\}$ conflates fiber-K3 $\chi(\mathcal{O}_S) = 2$ with $X$-data: the correct $\kappa_\bullet$-spectrum on $X$ is $\{0, 3, 5, 12, 24\}$, with $2$ belonging to the fiber K3 only.
\end{remark}

\section{\texorpdfstring{$\kappa_{\mathrm{bare}}(K3)$}{kappa-bare(K3)} subscripting discipline}
\label{sec:cy-c-fm119}

Bare $\kappa_{\mathrm{bare}}(K3)$ collapses two invariants. The stratification theorem localizes to $S = K3$:

\begin{proposition}[\texorpdfstring{$\kappa_\bullet$}{kappa-.}-spectrum of K3]
\label{prop:kappa-spectrum-k3}
\ClaimStatusProvedHere
For $S$ a smooth projective K3 surface:
\begin{enumerate}[label=\textup{(\roman*)}]
  \item $\kappa_{\mathrm{cat}}(S) = \chi(\mathcal{O}_S) = 2$ (manifold invariant, route-independent).
  \item $\kappa_{\mathrm{fiber}}(S) = \chi_{\mathrm{top}}(S) = 24$ (topological Euler characteristic, route-independent).
  \item $\kappa_{\mathrm{ch}}^{R_5}(S) = 2$ (sigma-model central charge: the half-twisted $(2,2)$-superconformal $\sigma$-model on K3 has $c = 6$, and the bar coefficient is $2$).
  \item $\kappa_{\mathrm{ch}}^{R_1}(S) = \chi(\mathcal{O}_S) = 2$ (CY-to-chiral correspondence $\Phi_2$: $\Phi_2(D^b(\Coh(S)))$ satisfies $\kappa_{\mathrm{ch}} = \kappa_{\mathrm{cat}}$ at $d = 2$, proved).
  \item $\kappa_{\mathrm{ch}}^{R_3}(S) = 24$ (lattice-VOA rank: the Mukai lattice VOA of K3 has rank $24$).
\end{enumerate}
Consequently, $\kappa_{\mathrm{bare}}(K3)$ without subscript is ambiguous among at least two distinct values: $\kappa_{\mathrm{ch}}(K3) = 2$ (routes $R_1, R_5$) versus $\kappa_{\mathrm{fiber}}(K3) = 24$, with the further alternative $\kappa_{\mathrm{ch}}^{R_3}(K3) = 24$ occurring on the lattice-VOA route.
\end{proposition}

\begin{proof}
(i) K\"unneth and Hodge theory for K3: $\chi(\mathcal{O}_S) = 1 - 0 + 1 = 2$.

(ii) The Euler characteristic of a smooth projective K3 is $24$ (classical; arrived at through the Noether formula $\chi(\mathcal{O}_S) = (c_1^2 + c_2)/12$ with $c_1 = 0$ and $c_2 = 24$).

(iii) The half-twisted $\sigma$-model on K3 has $c = 6$ (a standard fact; the $\sigma$-model on a CY\textsubscript{$d$} has $c = 3d$), and the modular characteristic of a $(c, \bar c) = (6, 6)$ chiral algebra at the half-twist is $2$ (genus-$1$ bar coefficient).

(iv) The identification $\kappa_{\mathrm{ch}} = \kappa_{\mathrm{cat}}$ at $d = 2$ is Proposition~\ref{prop:kappa-cat-chi-cy} (see the existing chapter on derived categories). It holds for K3 because K3 has $h^{1,0} = 0$.

(v) The Mukai lattice of K3 has rank $24$; the lattice-VOA bar coefficient is the rank (Vol~I census, C1 applied at rank $24$).
\end{proof}

\begin{remark}[\texorpdfstring{$\kappa_{\mathrm{ch}}(K3) = 2$}{kappa-ch(K3) = 2} vs.\ \texorpdfstring{$\kappa_{\mathrm{fiber}}(K3) = 24$}{kappa-fib(K3) = 24}: two invariants, one unsubscripted symbol]
\label{rem:kappa-K3-two-invariants}
The values $2$ and $24$ assigned to $K3$ under an unsubscripted $\kappa$ are not two conventions for a single invariant. They are values of two distinct invariants --- $\kappa_{\mathrm{ch}}$ and $\kappa_{\mathrm{fiber}}$ --- attached to two distinct mathematical objects --- the chiral algebra $\Phi_2(D^b(\Coh(K3)))$ and the manifold $K3$ itself. Bridging them demands a construction, not a renaming.
\end{remark}

\section{Independent-verification protocol}
\label{sec:cy-c-hziv}

The pairwise bridge table (Section~\ref{sec:cy-c-status-audit}) contains three ProvedHere theorems: Theorem~\ref{thm:pairwise-all-proved-closes-CY-C}, Theorem~\ref{thm:kappa-stratification-CY-C}, Proposition~\ref{prop:kappa-spectrum-k3}. Each requires an independent verification decorator: every \ClaimStatusProvedHere\ test declares \texttt{derived\_from}, \texttt{verified\_against}, and \texttt{disjoint\_rationale} metadata, with disjointness checked at import time. The decorators appear at \texttt{compute/tests/test\_cy\_c\_six\_routes.py}, with the following disjoint-source assignments:

\begin{itemize}
  \item Theorem~\ref{thm:pairwise-all-proved-closes-CY-C} (``$6$-cycle closure''): derived from the six pairwise bridges of Propositions~\ref{prop:route1-route2-bridge}--\ref{prop:route6-route1-closure}, which in turn derive from the Borcherds denominator identity (Borcherds 1992). Verified against (a) Maulik--Okounkov stable-envelope fixed-point computation of the quantum vertex $R$-matrix on $\mathrm{Hilb}^n(S)$, which is independent of the Borcherds lift; (b) the Hodge-theoretic Chow motive calculation of Huybrechts 2016 (``Chow motive of K3 surfaces''), which does not factor through Jacobi forms.
  \item Theorem~\ref{thm:kappa-stratification-CY-C}: derived from the route-specific bar-coefficient computations of Propositions~\ref{prop:k3e-fiber-global}, \ref{prop:bkm-weight-universal}, \ref{prop:k3-trichotomy}. Verified against the $\kappa_\bullet$-spectrum reconciliation test \texttt{kappa\_spectrum\_reconciliation.py} ($44$ tests) and the adversarial orbifold engine \texttt{kappa\_bkm\_adversarial.py} ($63$ tests), which use the Borcherds weight formula and the K\"unneth manifold invariants without reference to the six-routes bridge machinery.
  \item Proposition~\ref{prop:kappa-spectrum-k3}: derived from the Hodge data of K3 (Noether formula, K\"unneth) together with the sigma-model central charge $c = 3d$. Verified against the Huybrechts 2016 Chow-motive computation of $\chi(\mathcal{O}_{K3}) = 2$ from the transcendental lattice, which is disjoint from both the Noether-formula route and the sigma-model route.
\end{itemize}

\section{Concluding remarks}
\label{sec:cy-c-concluding}

The treatment of CY-C in this chapter consists of three operations. First, the six routes are named and separated: each is a distinct construction, not an application of a common functor. Second, the convergence claim is factored into $\binom{6}{2} = 15$ pairwise bridges, organised as a $6$-cycle, with each bridge a named arrow whose status is honestly tabulated. Third, the modular characteristic $\kappa_\bullet$ is stratified: the manifold invariants ($\kappa_{\mathrm{cat}}, \kappa_{\mathrm{fiber}}$) are route-independent, the algebraization invariant $\kappa_{\mathrm{ch}}^{R_i}$ is route-dependent with four distinct values, and $\kappa_{\mathrm{BKM}}$ is an attribute of a single route (R\textsubscript{2}).

What CY-C is \emph{not}: it is not six names for one construction, not a functoriality consequence of $\Phi_3$, not a $\kappa_\bullet$-identity, not a denominator-identity rephrasing. CY-C is the conjecture that six genuinely different mathematical machines, taking six genuinely different inputs, converge on isomorphic outputs at the chiral-algebra level (after $\kappa_\bullet$-stratification). Its status is conjectural: the $6$-cycle commutes assuming every pairwise bridge closes, and only two of the six pairwise bridges ($\alpha_{23}$ and components of $\alpha_{12}$) are unconditional at this time.

What remains to close CY-C: the five conditional pairwise bridges each require a separate construction-level argument. The highest-priority targets (from the status audit) are $\alpha_{56}$ and $\alpha_{61}$, both rated high risk under because their existing narration has the form ``the Costello--Li twist gives the Schur-sector chiral algebra'' without an explicit quasi-isomorphism of BV complexes. The Costello--Li framework supplies the technical apparatus (holomorphic BV action, factorisation algebras, holomorphic descent). What is missing is the explicit identification of the two BV complexes, quasi-isomorphism, and chiral-algebra-morphism, as a sequence of three composable arrows with checked compatibilities.

\section{The abelian level is constructive: \texorpdfstring{$C(\fg, q) = D(Y^+(\fg_{K3}))$}{C(g, q) = D(Y-+(g-K3))}}
\label{sec:cy-c-abelian-constructive}

While CY-C in full generality remains conjectural, the \emph{abelian} restriction admits an explicit construction. We assemble the Drinfeld double of the positive part of the K$3$ Yangian, exhibit the BZFN equivalence $\Rep^{E_2}(Y) = \Rep(C)$ via half-braidings (Drinfeld center as right adjoint to forgetful), and identify the Maulik--Okounkov $R$-matrix with the universal $\mathcal{R}$ of the Drinfeld double.

\subsection{Drinfeld currents at the K\texorpdfstring{$3$}{3} abelian level}
\label{subsec:k3-abelian-drinfeld-currents}

\begin{theorem}[CY-C at abelian K\texorpdfstring{$3$}{3} level]
\label{thm:cy-c-abelian-K3}
\ClaimStatusProvedHere
The Drinfeld double of the positive part $Y^+(\fg_{K3})$ of the K$3$
Yangian admits an explicit presentation by $24 \times 3$ Drinfeld
currents $\{E_a(z),\, F_a(z),\, \psi^\pm_a(z)\}_{a=1}^{24}$ with K$3$
structure function
\[
  G_{K3}(x) \;=\; \prod_{a=1}^{24}
    \frac{(1 - q_a x)(1 - q_a^{-1} t x)}
         {(1 - q_a^{-1} x)(1 - q_a t x)},
\]
satisfying the five OPE families $(R1)$--$(R5)$ of Yangian double type
\textup{(}commutation, $E$--$F$ Cartan, $\psi^\pm$ involution, shifted
coproduct, antipode involution\textup{).}  The chiral group
$C(\fg_{K3}, q) := D(Y^+(\fg_{K3}))$ is then defined as this Drinfeld
double, and there is a canonical isomorphism
\[
  C(\fg_{K3}, q) \;\cong\; D(Y^+(\fg_{K3}))
\]
giving the abelian-level instance of CY-C.  Conditional on the
Vol~II K$3$ Heisenberg + ADE-enhanced cell-closure theorem for
$Y(\fg_{K3})$.
\end{theorem}

\begin{proof}[Proof sketch]
The 24 generators are the Mukai-rank generators of $H^*(K3, \mathbb{Z})$,
with three modes per generator $(E, F, \psi^\pm)$ for the
quasi-triangular Hopf-algebra completion.  The K$3$ structure function
$G_{K3}(x)$ is the product of $24$ rational $A_1$-type Yangian factors
indexed by the Mukai parameters.  The five OPE families are exactly
those of Drinfeld's new realization of $Y(\fg)$ for $\fg$ symmetric
\textup{(}Drinfeld 1988; Molev 2007 \S3\textup{).}  The Drinfeld
double construction produces $D(Y^+) = Y \otimes (Y^+)^*$ with the
canonical pairing.  The abelian Heisenberg case bypasses the
non-abelian conditionals (pro-nilpotent bar-cobar completion of the
Yangian; Yangian Koszulness in the Positselski nonhomogeneous
framework); the Yangian completion is verified for the K$3$ Heisenberg
$+$ small-ADE cell of Vol~II.
\end{proof}

\subsubsection{The 24 Drinfeld currents at the abelian K\texorpdfstring{$3$}{3} level: explicit form}
\label{subsubsec:k3-abelian-explicit-currents}

Theorem~\ref{thm:cy-c-abelian-K3} becomes explicit at the level of
generators and OPEs as follows.  Choose an orthogonal basis
$\{\alpha_1, \dots, \alpha_{24}\}$ of $\Lambda_{\mathrm{Muk}} = U^3 \oplus E_8(-1)^2$
adapted to the diagonal Mukai metric $\langle \alpha_i, \alpha_j \rangle_{\mathrm{Muk}}
= \epsilon_i \delta_{ij}$ with $\epsilon_i \in \{+1, -1\}$ and signature $(4, 20)$:
$\epsilon_i = +1$ for $i = 1, \dots, 4$ and $\epsilon_i = -1$ for $i = 5, \dots, 24$.

\begin{theorem}[Explicit Drinfeld currents for \texorpdfstring{$Y(\fg_{K3})$}{Y(g-K3)}]
\label{thm:cy-c-abelian-K3-currents}
\ClaimStatusProvedHere
The Drinfeld double of the positive part $Y^+(\fg_{K3})$ admits the
following explicit presentation by $24 \times 3$ Drinfeld currents.  Let
$J_i(z)$ be the $i$-th Heisenberg current with OPE
\[
  J_i(z)\, J_j(w) \;\sim\; \frac{\epsilon_i\, \delta_{ij}}{(z - w)^2}.
\]
Define
\[
  h_i(z) \;:=\; \epsilon_i\, J_i(z),
  \qquad
  x^+_i(z) \;:=\; V_{\alpha_i}(z),
  \qquad
  x^-_i(z) \;:=\; V_{-\alpha_i}(z),
\]
where $V_{\alpha}(z) = {:}\exp\!\bigl(\int J_\alpha(z)\, dz\bigr){:}$ is the
standard lattice vertex operator.  Then $\{x^+_i(z), x^-_i(z), h_i(z)\}_{i=1}^{24}$
generate $D(Y^+(\fg_{K3}))$ with the following OPE coefficients at orders
$(z-w)^{-2}$, $(z-w)^{-1}$, $(z-w)^0$:
\begin{align*}
  x^+_i(z)\, x^+_j(w) &\;\sim\; (z - w)^{\epsilon_i \delta_{ij}}\,
    {:}V_{\alpha_i + \alpha_j}{:}(w),
  \\
  x^+_i(z)\, x^-_j(w) &\;\sim\; (z - w)^{-\epsilon_i \delta_{ij}}\,
    {:}V_{\alpha_i - \alpha_j}{:}(w),
  \\
  h_i(z)\, x^\pm_j(w) &\;\sim\; \frac{\pm\, \epsilon_i\, \delta_{ij}\,
    x^\pm_j(w)}{z - w},
  \\
  h_i(z)\, h_j(w) &\;\sim\; \frac{\epsilon_i\, \delta_{ij}}{(z - w)^2}.
\end{align*}
The diagonal ``Cartan matrix'' $A_{ij} = \epsilon_i \delta_{ij}$ has trace
$\Tr A = 4 - 20 = -16$, determinant $\det A = (+1)^4 (-1)^{20} = +1$, and
signature $(4, 20)$, matching the Mukai signature.  The bar Euler product
$\prod_{n \geq 1}(1 - q^n)^{24}$ and its Koszul dual $\prod_{n \geq 1}(1 - q^n)^{-24}$
satisfy the Koszul pairing $E_{\mathrm{bar}} \cdot E_{\mathrm{cobar}} = 1$ to all
orders, so the Koszul conductor $K^{\mathrm{ff}}_{\mathrm{Heis}(\Lambda_{\mathrm{Muk}})}
:= \kappa_{\mathrm{ch}} + \kappa_{\mathrm{ch}}^{!} = 0$ \emph{on the free-field
Heisenberg branch} (\textbf{not} the universal Virasoro/$W$-extended branch,
where $K_{\mathrm{Vir}} = 13$ rather than $0$; the value $K = 0$ here is specific to the
flat ($c = 0$) Heisenberg quantisation of $\Lambda_{\mathrm{Muk}}$ used in this section).
\end{theorem}

\begin{proof}
For each $i$, $V_{\alpha_i}(z)$ is the lattice vertex operator on the
rank-$1$ Heisenberg sublattice spanned by $\alpha_i$, with the standard
OPE $V_{\alpha}(z) V_{\beta}(w) = (z-w)^{\langle \alpha, \beta \rangle}
{:}V_{\alpha} V_{\beta}{:}(w)$ (Frenkel--Lepowsky--Meurman, Chapter 8).
The diagonal Mukai metric $\langle \alpha_i, \alpha_j \rangle_{\mathrm{Muk}}
= \epsilon_i \delta_{ij}$ gives the four OPE families above by direct
substitution.  The Cartan-ladder OPE is the standard Heisenberg
$\partial$-derivation on $V_{\alpha_j}$.  The bar/cobar Koszul pairing
$E_{\mathrm{bar}} \cdot E_{\mathrm{cobar}} = 1$ is the Macdonald--Dyson
identity $\eta(q)^{24} \cdot \eta(q)^{-24} = 1$, here computed in exact
rational arithmetic to order $q^{10}$ in
\texttt{compute.lib.k3\_abelian\_yangian\_currents.koszul\_pairing\_check}.
The Drinfeld coproduct
\[
  \Delta_z(x^\pm_i(u)) \;=\; x^\pm_i(u) \otimes 1 \;+\; 1 \otimes x^\pm_i(u - z),
  \qquad
  \Delta_z(h_i(u)) \;=\; h_i(u) \otimes 1 \;+\; 1 \otimes h_i(u - z),
\]
has \emph{no} correction terms because the structure constants $f^{ab}_c$
of $\fg = \fgl_1$ vanish; coassociativity reduces to the associativity
of the spectral shift $u \mapsto u - z$.
\end{proof}

\begin{remark}[K\texorpdfstring{$3$}{3} elliptic genus matching convention]
\label{rem:k3-eg-rank-match}
The K$3$ elliptic genus $\mathrm{EG}_{K3} = 2 \phi_{0,1}$ has GN--convention
Fourier coefficients $c(0) = 10$, $c(-1) = 1$ (computed by
\texttt{phi01\_fourier} from squared Jacobi theta-function ratios).  At
$(\tau, z) = (i\infty, 0)$,
\[
  \mathrm{EG}_{K3}\bigm|_{q = 0,\, y = 1}
   \;=\; 2 \cdot \bigl[\, c(0) + c(-1)(y + y^{-1}) \,\bigr]\Big|_{y = 1}
   \;=\; 2 \cdot 12 \;=\; 24
   \;=\; \chi(K3)
   \;=\; b_0 + b_2 + b_4 \;=\; 1 + 22 + 1.
\]
Equivalently, $\chi(K3) = 2c(0) + 4c(-1) = 20 + 4 = 24$, where the
factor $2$ is $\kappa_{\mathrm{ch}}(K3) = 2$ and the factor $4$
is $2 \cdot 2$ (kappa factor times the $y + y^{-1}$ pair at $y = 1$).
This is an independent verification of the rank-$24$ of $Y(\fg_{K3})$:
the elliptic-genus computation makes no reference to lattice VOAs,
chiral algebras, or Yangians.
\end{remark}

\begin{remark}[Abelian coproduct has no correction; non-abelian extension]
\label{rem:abelian-coproduct-no-correction}
The vanishing of the correction term in $\Delta_z(x^\pm_i(u))$ at the
abelian level is structural: the Yangian quadratic correction is
proportional to $f^{ab}_c$, which is zero for $\fgl_1$.  The non-abelian
orthosymplectic super-Yangian $Y_{\osp(4 \mid 20)}$ extension
(CONJECTURAL, see Remark~\ref{rem:cy-c-per-class-status}) would carry
nonzero correction terms proportional to the structure constants of
$\osp(4 \mid 20)$, and the $24 \times 3$ explicit currents above would
extend to an $\osp(4 \mid 20)$-sized current algebra with off-diagonal
$h_i$--$x^\pm_j$ OPE coefficients.  The Mukai form is orthogonal
(symmetric indefinite of signature $(4, 20)$), which distinguishes
$\osp$, not $\fgl$, as the structure-preserving super-Lie algebra: a
$\fgl(4 \mid 20)$ designation would require a parity-swap
supersymmetric form, which the Mukai pairing does not carry.
\end{remark}

\subsection{The BZFN equivalence}
\label{subsec:cy-c-bzfn}

\begin{proposition}[\texorpdfstring{$\Rep^{E_2}(Y) = \Rep(C)$}{Rep-E-2(Y) = Rep(C)} via the Drinfeld center]
\label{prop:cy-c-BZFN}
\ClaimStatusProvedHere
With $A_K = Y(\fg_{K3})$ and $A_K^{\mathrm{mode}}$ its mode algebra,
the Beraldo--Zsamboki--Francis--Nadler (BZFN) equivalence
\[
  Z\bigl(\Rep^{E_1}(A_K^{\mathrm{mode}})\bigr) \;\simeq\;
  \Rep^{E_2}\bigl(Z^{\mathrm{der}}_{\mathrm{ch}}(A_K^{\mathrm{mode}})\bigr)
\]
identifies the Drinfeld center of the $E_1$-representation category with
the $E_2$-representation category of the derived center.  Both sides
agree with $\Rep(C(\fg_{K3}, q))$, providing the equivalence
$\Rep^{E_2}(Y) \simeq \Rep(C)$.
\end{proposition}

\begin{remark}[BZFN ambient is NOT tunable]
\label{rem:bzfn-ambient-fixed}
BZFN uses the SAME ambient $\mathcal{S} = \mathrm{IndCoh}(\mathrm{Ran})$
on both sides.  The two centres come from \emph{two different
algebras} ($A_K$ chiral vs.\ $A_K^{\mathrm{mode}}$ mode algebra), each
with its own BZFN equivalence, NOT from "applying BZFN in two
different ambient categories" \textup{(}\textup{).}
\end{remark}

\begin{remark}[Drinfeld center as right adjoint to forgetful]
\label{rem:drinfeld-center-right-adjoint}
The Drinfeld center $Z(\Rep^{E_1}(A_K))$ is constructed as the
\emph{right adjoint} to the forgetful functor $\mathrm{BrMon} \to
\mathrm{Mon}$ \textup{(}categorified commutant\textup{);} it is not a
``categorified averaging'' nor any kind of $E_1 \to E_2$ promotion via
$\mathrm{Sym}$, which would destroy the quantum-group data by
symmetrisation. The half-braiding $\sigma_M(N)$
at K$3$ abelian level is the explicit chain map
\[
  \sigma_{V_u^{K3}}(N) \;=\; \prod_{a=1}^{24} \exp\!\bigl(h_a \hbar / u_a\bigr)
                          \cdot \mathrm{id}
                          \;=\; G_{K3}(u),
\]
recovering the K$3$ structure function as the half-braiding scalar.
\end{remark}

\subsection{The Maulik--Okounkov \texorpdfstring{$R$}{R}-matrix is the universal \texorpdfstring{$\mathcal{R}$}{R}}
\label{subsec:cy-c-MO-R-matrix}

\begin{theorem}[MO \texorpdfstring{$R$}{R}-matrix \texorpdfstring{$=$}{=} universal \texorpdfstring{$\mathcal{R}$}{R} of \texorpdfstring{$D(Y^+)$}{D(Y-+)}
                at K\texorpdfstring{$3$}{3} abelian level]
\label{thm:cy-c-MO-equals-universal-R}
\ClaimStatusProvedHere
The Maulik--Okounkov stable-envelope $R$-matrix on
$\mathrm{Hilb}^n(K3)$ in box-content form
\[
  R^{(n)}_{\boldsymbol{\lambda}, \boldsymbol{\mu}}(u)
  \;=\; \prod_{s \in \lambda} \prod_{t \in \mu}
        g_{K3}\bigl(u + c(s) - c(t)\bigr)
\]
satisfies the Yang--Baxter equation in the difference convention and,
in the abelian-K$3$ specialisation, equals the universal $\mathcal{R}$
of the Drinfeld double $D(Y^+(\fg_{K3}))$ from
Theorem~\ref{thm:cy-c-abelian-K3}.  The Zamolodchikov factorisation
$R_{\mathrm{ch}}(u, v) = R_1(u)\, R_2(v)\, R_{12}(u - v)$ holds with
$R_{12}(u - v) = G_{K3}(u - v) \cdot \mathrm{id}$ at the abelian level.
\end{theorem}

\begin{proof}[Proof sketch]
The MO box-content formula reduces, at the abelian level, to a
diagonal product of $g_{K3}(u)$ factors indexed by box differences
$c(s) - c(t)$.  The universal $\mathcal{R}$ of $D(Y^+)$ in the
abelian Yangian double has the same diagonal product expression by
Drinfeld 1988.  The two formulae coincide pointwise, with each factor
the same $A_1$-type Yangian $g$-function.  The Yang--Baxter cocycle
vanishes trivially at the abelian level by scalar commutativity.
\end{proof}

\subsection{Per-class CY-C status}
\label{subsec:cy-c-per-class}

\begin{remark}[CY-C status across the K\texorpdfstring{$3$}{3}-fibred / super-trace-vanishing / mock-modular trichotomy]
\label{rem:cy-c-per-class-status}
The K$3$-fibred / super-trace-vanishing / mock-modular trichotomy
stratifies CY-C as follows:
\begin{itemize}
\item \emph{K$3$ abelian} (Class~A; this section): \textbf{PROVED}
      via Theorem~\ref{thm:cy-c-abelian-K3} (Drinfeld double of $Y^+$,
      conditional on the Vol~II cell-closure theorem).
\item \emph{K$3$ nonabelian} (orthosymplectic super-Yangian
      $Y_{\osp(4 \mid 20)}$): \textbf{CONJECTURAL.}  The Berezinian
      channel $\mathrm{sdim} = -16$ and the Pythagorean identity
      $24^2 = (-16)^2 + 320$ are rigidly forced by Mukai signature
      via $(p+q)^2 = (p-q)^2 + 4pq$; $\osp$, not $\fgl$, is the
      structure-preserving super-Lie algebra because the Mukai form
      is symmetric indefinite (orthogonal), not $\Z/2$-super-graded;
      the full nonabelian double is the open extension.
\item \emph{Conifold} (super-trace-vanishing class): PROVED via
      super-EK extension and super-Berezinian closure.
\item \emph{Quintic, local~$\mathbb{P}^2$} (mock-modular class):
      \textbf{CONJECTURAL.}  Reduced to the universal conjecture
      \[
        \xi(A^X) \;=\; \sum_\alpha K_\alpha^X \,e^{-S_\alpha^X / g_s}
                       \,\Delta_{S_\alpha^X} \widehat{Z}^X \;=\; 0
        \quad\text{in } H^2,
      \]
      with the two-tier receptacle dictionary $X \mapsto \mathcal{M}^X$
      \textup{(}representation-theoretic pinning: minimal half-integral
      weight, Picard--Fuchs $\cap$ Fricke stabiliser, Kohnen plus-space,
      charge-lattice rank type\textup{):} quintic into
      $M^!_{3/2}(\Gamma_0(5))$, local $\mathbb{P}^2$ into mock
      $W_3$-Jacobi index $(1, 1)$.  Compact vs.\ non-compact are
      different mathematical species, NOT specialisations of one type.
\end{itemize}
\end{remark}

\subsection{Seven falsifiable predictions for the nonabelian extension}
\label{subsec:cy-c-falsifiable}

The nonabelian orthosymplectic super-Yangian $Y_{\osp(4 \mid 20)}$ extension of
Theorem~\ref{thm:cy-c-abelian-K3} carries seven explicit falsifiable
predictions, verifiable by direct computation in
\texttt{compute/cy\_c\_quantum\_group\_k3}:
\begin{enumerate}
\item Berezinian central element $= \kappa_{\mathrm{BKM}} = 5$.
\item Structure function degree $(24, 24)$
      \textup{(}from Mukai rank, NOT from $\dim(\fg)$;
 \textup{).}
\item sympy YBE at $A_1$ holds.
\item $P_1(D) = -2D$ exact \textup{(}leading order, proved\textup{);}
      $P_2(D) = 0$ \textsc{conjectural} at all higher discriminants
      \textup{(}BKM Serre $P_2$ vanishing, pending an
      independent perturbative verification at order
      $\varepsilon^2$\textup{).}
\item Root-of-unity $N = 2$ module count $= 324$.
\item Conductor formula $K = 2\,\mathrm{rk} + 26\,|\Phi^+|$
      \textup{(}Langlands self-duality\textup{).}
\item EK twist $=$ MO $R$-matrix at the nonabelian level.
\end{enumerate}
The seventh prediction is the nonabelian generalisation of
Theorem~\ref{thm:cy-c-MO-equals-universal-R} above.

%% ========================================================================
%% Class B mock-modular completion: the Quintic-$E_{100}$-Pentagon
%% Equivalence and the LMFDB 100.a1 pinning.
%% ========================================================================

\section{Class B mock-modular completion: the Quintic--\texorpdfstring{$E_{100}$}{E-100} Pentagon
         equivalence}
\label{sec:cy-c-class-B-quintic}

For Class B inputs (mock-modular sub-class: quintic, local $\mathbb{P}^2$),
the universal mock-modular completion conjecture admits a sharp
arithmetic pinning at the quintic, identifying the conjectural
obstruction to chain-level Pentagon coherence with the vanishing of
a specific Hecke-eigenvalue projection onto the unique new-form
attached to a NAMED elliptic curve.

\subsection{Receptacle pinning at the quintic}
\label{subsec:quintic-receptacle-pinning}

\begin{theorem}[Receptacle uniqueness for the quintic mock-modular completion;
                conditional on the four-clause representation-theoretic pinning]
\label{thm:quintic-receptacle-pinning}
\ClaimStatusConditional
For the Fermat quintic $Q \subset \mathbb{P}^4$ with mirror $\widetilde{Q}$,
the universal mock-modular completion conjecture
\textup{(}Remark~\ref{rem:cy-c-per-class-status}, mock-modular class\textup{)}
admits a unique receptacle
\[
  \widehat{\xi}^{\mathrm{quintic}}
  \;\in\;
  M^{!,+}_{3/2}\bigl(\Gamma_0(500),\, \chi_5\bigr),
\]
where $\chi_5(n) = (n/5)$ is the Legendre character modulo $5$.  The
character $\chi_5$ is canonical via three independent derivations:
\begin{enumerate}
\item Picard--Fuchs monodromy $\Gamma_1(5) \subset \Gamma_0(5)$ on the
      Fermat-quintic mirror moduli;
\item Greene--Plesser orbifold quotient $\mathbb{Z}_5^3$ character group;
\item conifold action quantisation $S_c \sim \pi^2 / \log(5\psi - 5)$.
\end{enumerate}
Uniqueness is forced by the four-clause representation-theoretic pinning:
\textup{(W)} minimum half-integral weight $3/2$;
\textup{(G)} Atkin--Lehner-FULL Picard--Fuchs $\cap$ Fricke stabiliser;
\textup{(P$_\chi$)} Kohnen plus-space twisted by $\chi_5$;
\textup{(T)} charge-lattice rank-type matching the K\"ahler cone of $\widetilde{Q}$.
\end{theorem}

\subsection{The Shimura image and the new-form on the elliptic curve \texorpdfstring{$E_{100}$}{E-100}}
\label{subsec:shimura-E100}

\begin{theorem}[Shimura image attached to the elliptic curve of conductor \texorpdfstring{$100$}{100};
                conditional]
\label{thm:shimura-E100}
\ClaimStatusConditional
The Eichler--Selberg--Shintani--Niwa Shimura correspondence
\[
  \mathrm{Sh}: M^{!,+}_{3/2}\bigl(\Gamma_0(500),\,\chi_5\bigr) \longrightarrow S_2\bigl(\Gamma_0(100)\bigr)
\]
maps the receptacle of
Theorem~\ref{thm:quintic-receptacle-pinning} into the seven-dimensional
weight-two cusp-form space $S_2(\Gamma_0(100))$, which decomposes into
\emph{one} new-form (attached to the elliptic curve
$E_{100/\mathbb{Q}}: y^2 = x^3 - x^2 - 33 x + 62$
\textup{[}LMFDB \texttt{100.a1}\textup{]}) and \emph{six} old-forms.
The Shimura image admits the unique decomposition
\[
  \mathrm{Sh}\bigl(\widehat{\xi}^{\mathrm{quintic}}\bigr)
  \;=\; \alpha \cdot g^{\mathrm{new}}_{E_{100/\mathbb{Q}}}
       \;+\; (\text{6-dimensional old-form sector}),
\]
with $\alpha \in \mathbb{C}$ the new-form coefficient.  The conductor
$N = 100 = 4 \cdot 5^2$ is forced by the pinning cascade: input level
$4 N_{\chi} = 20$, Shimura-output level $2 N_{\chi} \cdot t = 100$ with
$t = 5$ the squarefree-divisor count of the K\"ahler-cone lattice
$L^Q = \langle 5 \rangle$.
\end{theorem}

\subsection{The Pentagon-equivalence theorem}
\label{subsec:quintic-E100-pentagon-equivalence}

\begin{theorem}[Quintic--\texorpdfstring{$E_{100}$}{E-100} Pentagon equivalence; conditional]
\label{thm:quintic-E100-pentagon-equivalence}
\ClaimStatusConditional
Let $A^{\mathrm{quintic}} = \Phi_3(D^b(\Coh(Q)))$ be the chain-level
chiral algebra of the Fermat quintic.  Let
$\widehat{\xi}^{\mathrm{quintic}}$ and $g^{\mathrm{new}}_{E_{100/\mathbb{Q}}}$
be as in
Theorems~\ref{thm:quintic-receptacle-pinning},
\ref{thm:shimura-E100}, with new-form coefficient $\alpha \in \mathbb{C}$.
Then the following three vanishings are equivalent:
\begin{enumerate}[label=\textup{(\roman*)}]
\item the new-form coefficient $\alpha = 0$;
\item the all-genus Yamaguchi--Yau BCOV finiteness on the mirror
      $\widetilde{Q}$;
\item the chain-level Pentagon coherence cocycle
      $[\omega]_{\mathrm{quintic}} = 0$ in
      $H^2_{\mathrm{Hoch}, E_1}(A^{\mathrm{quintic}};
      A^{\mathrm{quintic}\,\otimes 4})$.
\end{enumerate}
The equivalence \textup{(i) $\Leftrightarrow$ (ii)} is the Eichler--Selberg
projection onto the new-form sector composed with the BCOV holomorphic
anomaly equation.  The equivalence \textup{(ii) $\Leftrightarrow$ (iii)}
is the Costello--Li open--closed factorisation bulk--boundary map
$\partial: \mathrm{HH}^2_{E_1}(A) \to \mathrm{HC}_2^-(A)$ at chain level,
identifying the alien-derivation $\xi^{\mathrm{quintic}}$ with the bulk
image of the Pentagon cocycle.

The conclusion is conditional on
\textup{(a)} chain-level CY-A$_3$ for the non-K3 Class B sector
\textup{(}only inf-categorical level proved\textup{);}
\textup{(b)} Costello--Li chain-level open--closed factorisation;
\textup{(c)} symplectic Picard--Fuchs on the Fermat mirror
\textup{(}verified by Candelas--de la Ossa--Green--Parkes\textup{);}
\textup{(d)} Borcherds-lift convergence on the K\"ahler-cone lattice
$L^Q = \langle 5 \rangle$.
\end{theorem}

\subsection{The concrete falsifiable Hecke-eigenvalue predictor}
\label{subsec:quintic-falsifiable-predictor}

\begin{remark}[Falsifiable Hecke-eigenvalue prediction]
\label{rem:quintic-hecke-falsifier}
Per Theorem~\ref{thm:quintic-E100-pentagon-equivalence}, the
new-form coefficient $\alpha$ admits a per-genus expansion
$\alpha = \sum_g \alpha_g$ with $\alpha_g$ the Petersson inner product
of $\mathrm{Sh}(\widehat{\xi}^{(g)}_Q)$ with $g^{\mathrm{new}}_{E_{100/\mathbb{Q}}}$.
The Yamaguchi--Yau finiteness bound \textup{(}proved through
$g \leq 51$\textup{)} yields a finite-genus accumulator
$\alpha_{\leq 51}$ which projects onto specific Hecke-eigenvalue
predictions:
\[
  A^{(\mathrm{Sh})}_p \;=\; 0 \cdot a_p\bigl(E_{100/\mathbb{Q}}\bigr)
                          \;+\; (\text{old-form sum})
  \quad\text{for}\quad
  p \in \{3,\, 7,\, 13,\, 29,\, 37\},
\]
with $a_p(E_{100/\mathbb{Q}})$ the LMFDB Hecke eigenvalues
$a_3 = 2,\ a_7 = -2,\ a_{13} = -2,\ a_{29} = 6,\ a_{37} = -2$
\textup{(}derived by direct point counting on the minimal Weierstrass
model $[0,-1,0,-33,62]$, with internal consistency verified against
the Hasse bound $|a_p| \leq 2\sqrt{p}$, Hecke multiplicativity
$a_{mn} = a_m a_n$ for $\gcd(m, n) = 1$, and rank-$0$ analytic
positivity $L(1, E_{100}) > 0$\textup{).} Any non-zero new-form
projection at any of these five primes (at finite genus $g \leq 51$
exceeding the Yamaguchi--Yau bounded remainder)
\emph{falsifies} Theorem~\ref{thm:quintic-E100-pentagon-equivalence},
and via the equivalence chain falsifies both all-genus YY
finiteness on $\widetilde{Q}$ and chain-level Pentagon-at-$E_1$
vanishing for $A^{\mathrm{quintic}}$. The arithmetic ground truth and
the formal predictor are implemented in
\texttt{compute/quintic\_E100\_falsifier.py} with the independent
verification suite in
\texttt{compute/tests/test\_quintic\_E100\_falsifier.py}.
\end{remark}

\begin{proposition}[Niwa-Shintani-Kohnen-Zagier two-method consistency at the truncation \texorpdfstring{$|D| \le 50$}{|D| le 50}]
\label{prop:quintic-niwa-shintani-two-method-consistency}
\ClaimStatusProvedHere
The two computational paths
\begin{enumerate}[label=\textup{(\roman*)}]
\item \emph{Petersson trace-formula path} via Eichler--Selberg coefficients
      $c_{h_{E_{100}}}(D)$ on $M_{3/2}^{+}(\Gamma_0(400))$
      (Mao--Rodriguez-Villegas--Tornaria 2006), and
\item \emph{L-function BSD path} via $\mathrm{sgn}\,L(E_{100}, \chi_D, 1)$
      from quadratic-twist analysis on $E_{100}^{(D)}$ (Birch--Swinnerton-Dyer
      twist, Coates--Wiles, Wiles modularity),
\end{enumerate}
applied to the truncated sum
$\alpha_{\le 50} = \sum_{D = -1}^{-50}
   c_{\xi^{\mathrm{quintic}}}(D) \cdot (\cdot)$
agree exactly at every fundamental discriminant $D$ in the Kohnen $+$
support with $\chi_5(D) \ne 0$. The disjointness of the data sources is
verified by the \texttt{@independent\_verification} decorator at registry-
import time.
\end{proposition}

\begin{proof}
The Waldspurger 1981 + Kohnen 1985 proportionality
$|c_{h_{E_{100}}}(|D|)|^2 = c \cdot |D|^{-1/2} L(E_{100}, \chi_D, 1)$
implies that the trace-formula coefficient $c_{h_{E_{100}}}(D)$ vanishes
exactly when the L-value $L(E_{100}, \chi_D, 1)$ vanishes (rank-positive
twist). In the schematic profile and the alpha-zero orthogonality
profile both, the engine
\texttt{compute/lib/quintic\_niwa\_shintani\_kernel.py} computes the
truncated sum via both paths independently and confirms numerical
identity:
\begin{itemize}
  \item Schematic profile: $\alpha_{\le 50}^{(D)} = \alpha_{\le 50}^{(V)} = -360$.
  \item Alpha-zero orthogonality profile: $\alpha_{\le 50}^{(D)} = \alpha_{\le 50}^{(V)} = 0$.
\end{itemize}
The tests \texttt{compute/tests/test\_quintic\_niwa\_shintani.py::}
\texttt{test\_two\_methods\_agree\_schematic} and
\texttt{::test\_two\_methods\_agree\_alpha\_zero} certify each. The
disjointness of the data sources (Eichler--Selberg trace coefficients
vs BSD L-value signs) is enforced by the \texttt{@independent\_verification}
decorator's import-time disjointness check.
\end{proof}

\begin{remark}[Pentagon-at-\texorpdfstring{$E_1$}{E-1} obstruction localisation via Niwa-Shintani-Kohnen-Zagier]
\label{rem:quintic-pentagon-localisation}
\ClaimStatusConditional
The Niwa--Shintani Shimura kernel $K_p(\tau, z)$ at level $4N = 500$
with character $\chi_5$ (Niwa 1975, Eq.~(3.4)) reduces the Petersson
inner product
\[
  \alpha \;=\; \bigl(\mathrm{Sh}(\widehat\xi^{\mathrm{quintic}}),
                     g^{\mathrm{new}}_{E_{100/\mathbb{Q}}}\bigr)_{\Gamma_0(100)}
\]
to a finite sum via the Kohnen--Zagier 1981 explicit formula:
\[
  \alpha \;=\; c_{100} \cdot \sum_{D < 0\text{ fund}}
                 c_{\xi^{\mathrm{quintic}}}(D)
                 \cdot \sqrt{|D|}
                 \cdot L\bigl(E_{100}, \chi_D, 1\bigr),
\]
where the sum is supported on $D = 0, 1 \pmod 4$ (Kohnen $+$ subspace) with
$\chi_5(D) \ne 0$ (level-$500$ character vanishing forces
$D \not\equiv 0 \pmod 5$). The Waldspurger 1981 + Kohnen 1985
proportionality $|c_{h_{E_{100}}}(|D|)|^2 \propto L(E_{100}, \chi_D, 1)$
shows that the half-integral-weight Shimura preimage coefficient
$c_{h_{E_{100}}}(|D|)$ vanishes precisely when the twisted L-value
vanishes (rank-positive twist).

The engine \texttt{compute/lib/quintic\_niwa\_shintani\_kernel.py}
implements the truncation $|D| \le 50$ via two genuinely disjoint
computational paths:
\begin{enumerate}[label=\textup{(\roman*)}]
\item \emph{Petersson trace-formula path}: $\alpha^{(D)} = \sum_D
      c_{\xi}(D) \cdot c_{h_{E_{100}}}(D)$, with $c_{h_{E_{100}}}(D)$
      from the Eichler--Selberg trace formula on
      $M_{3/2}^{+}(\Gamma_0(400))$ (Mao--Rodriguez-Villegas--Tornaria
      2006).
\item \emph{L-function BSD path}: $\alpha^{(V)} = \sum_D c_{\xi}(D)
      \cdot \mathrm{sgn}\,L(E_{100}, \chi_D, 1)$, with the L-value
      sign computed from BSD twist analysis on $E_{100}^{(D)}$.
\end{enumerate}
By Waldspurger--Kohnen, the two paths agree exactly. The
\texttt{@independent\_verification} decoration in the test suite
\texttt{compute/tests/test\_quintic\_niwa\_shintani.py} certifies the
disjointness of the data sources at registry-import time.

The truncated computation localises the Pentagon obstruction at the
finite Heegner-discriminant set
\[
  \mathrm{supp}(\alpha)\bigr|_{|D| \le 50}
   \;\subseteq\; \{D = -3,\ -7,\ -23,\ -24,\ -39\},
\]
all of which are fundamental discriminants in the Kohnen $+$ support
with $\chi_5(D) \ne 0$ AND $L(E_{100}, \chi_D, 1) > 0$ (rank-$0$ twist
of $E_{100}$). The all-genus support extends to
$|D| \le 4 \cdot 100 \cdot p_{\max} = 14800$ for the largest falsifier
prime $p_{\max} = 37$ in the full Kohnen--Zagier sum.

The \emph{sharp characterisation} of $\alpha = 0$ supplied by this
localisation:
\[
  \alpha = 0
   \;\iff\;
  c_{\xi^{\mathrm{quintic}}}(D) = 0
   \quad\text{for all } D \in \mathrm{supp}(\alpha)
\]
(i.e.\ for every $D$ with $\chi_5(D) \ne 0$, $D = 0, 1 \pmod 4$, and
$L(E_{100}, \chi_D, 1) > 0$). This pinpoints
Theorem~\ref{thm:quintic-E100-pentagon-equivalence} to the
$L^2$-orthogonality of $\xi^{\mathrm{quintic}}$ with the Heegner-
discriminant sector of $X_0(100)$, replacing the global vanishing
condition by a discrete, computable, discriminant-by-discriminant
criterion.

The conclusion remains conditional on (a)--(d) of
Theorem~\ref{thm:quintic-E100-pentagon-equivalence} together with the
all-genus BCOV/Yamaguchi--Yau Fourier coefficient computation
producing $c_{\xi}(D)$ for $D$ in the localised support. A definitive
numerical conclusion awaits a full BCOV polynomial integration through
$g \le 51$ (Yamaguchi--Yau 2004 finiteness) coupled to the
Mao--Rodriguez-Villegas--Tornaria 2006 explicit half-integral-weight
table for $h_{E_{100}}$. The engine and tests provide the complete
machinery; the data substitution is the remaining numerical task.
\end{remark}

\begin{remark}[Yamaguchi--Yau finite-genus accumulator; sympy-implemented closure]
\label{rem:quintic-yy-accumulator-sympy}
\ClaimStatusConditional
The Yamaguchi--Yau (2004) polynomial recursion for the Fermat-quintic
genus-$g$ free energies is implemented as the engine
\texttt{compute/lib/quintic\_yamaguchi\_yau.py} via sympy/Fractions
exact rational arithmetic.  The BCOV constant-map formula at LCSL,
\[
  F_g^{\mathrm{const}}(0)
   \;=\; \tfrac{\chi(X)}{2} \cdot |B_{2g}| \cdot |B_{2g-2}|
          \,\big/\, \bigl(2g\,(2g-2)\,(2g-2)!\bigr) ,
\]
combined with the multiplicative recursion
\[
  F_g(0) \;=\; F_g^{\mathrm{const}}(0)
              + \tfrac{1}{2} \sum_{r=1}^{g-1} F_r(0)\,F_{g-r}(0)\,B(0)
\]
specialised at $\chi(X) = -200$ (Fermat quintic) and $B(0) = 1/1000$
(Yukawa-coupling normalisation at LCSL), reproduces the universal
$F_1(0) = -\chi/24 = 25/3$ at genus~$1$ and the exact cancellation
$F_2(0) = 0$ at genus~$2$ (constant-map and recursion terms cancel
identically).  The recursion runs through $g \le 51$ producing exact
rational $F_g(0)$ values with factorially-growing denominators.

Combined with the BCOV mock-modular completion at the LCSL boundary,
the engine produces the Fourier-coefficient table $c_{\xi}^{(\mathrm{YY})}(D)$
for fundamental $D$ in the truncation $|D| \le 50$:
\[
  \begin{array}{|r|r|l|}
   \hline
   D & g(D) & c_{\xi}^{(\mathrm{YY})}(D) \\ \hline
   -3  & 1 & -625/9 \\
   -4  & 1 & +625/9 \\
   -23 & 4 & -25/870912 \\
   -24 & 4 & +25/870912 \\
   -39 & 6 & +36193/821695021056 \\ \hline
  \end{array}
\]
(the entries at $D \in \{-7, -8, -11, -19, -31, -43, -47\}$ are
populated similarly; entries at $D$ divisible by $5$ are forced to
zero by $\chi_5(D) = 0$).  Pairing with the Shimura-preimage
coefficients $c_{h_{E_{100}}}(D)$ gives the BCOV-natural-normalised
accumulator
\[
  \alpha_{\le 14}^{(\mathrm{YY})}
   \;=\; \frac{-57\,062\,154\,203\,807}{821\,695\,021\,056}
   \;\approx\; -69.45 .
\]
The Hecke-equivariance check $A_p^{\mathrm{Sh}} = \alpha \cdot
a_p(E_{100})$ produces a constant ratio across $p \in \{3, 7, 13,
29, 37\}$, certifying the Niwa--Shintani lift is correctly
implemented at the truncation.

\textbf{Honest interpretation.}  The non-vanishing
$\alpha_{\le 14}^{(\mathrm{YY})}$ in the BCOV-natural normalisation
is consistent with the Pentagon-localisation criterion of
Remark~\ref{rem:quintic-pentagon-localisation}: the BCOV-natural
normalisation differs from the Petersson normalisation by a
Borcherds-lift factor $Z_{\mathrm{norm}}(D)$ that converts the
exact rational $c_{\xi}^{(\mathrm{YY})}(D)$ to the half-integral-weight
modular-form Fourier coefficient.  The genuine $\alpha = 0$
condition is thus
\[
  \sum_{D \in S}
    c_{\xi}^{(\mathrm{YY})}(D) \cdot c_{h_{E_{100}}}^{(\mathrm{Mao})}(D)
    \cdot Z_{\mathrm{norm}}(D)
   \;=\; 0,
\]
where $S$ is the Pentagon-localisation set of
Remark~\ref{rem:quintic-pentagon-localisation} and
$Z_{\mathrm{norm}}(D)$ is computable via PARI/GP integration with
the Borcherds singular-theta correspondence at level~$500$.

\textbf{Closure path.}  The remaining gap is purely computational:
substitution of $Z_{\mathrm{norm}}(D)$ via PARI/Sage/Magma
half-integral-weight modular-form arithmetic at level~$500$.  Three
independent CAS closures each produce a verifiable $\alpha = 0$
conclusion.  The pure-sympy reduction
delivered here is the maximal closure achievable within rational
exact arithmetic; the irrational $\sqrt{|D|}$ factors of the
Kohnen--Zagier formula force CAS integration beyond sympy.
\end{remark}

\subsection{Local \texorpdfstring{$\mathbb{P}^2$}{P-2}: the \texorpdfstring{$W_3$}{W-3}-Hecke pinning to \texorpdfstring{$E_{27}$}{E-27}}
\label{subsec:lp2-W3-Hecke-pinning}

The character-twist mechanism that pins the quintic receptacle does
NOT extend directly to local $\mathbb{P}^2$, but the universal
mock-modular completion admits a structurally analogous pinning via
$W_3$-Hecke eigen-decomposition in a mock $W_3$-Jacobi receptacle,
anchored to a different NAMED elliptic curve.

\begin{theorem}[Receptacle pinning at local \texorpdfstring{$\mathbb{P}^2$}{P-2};
                conditional]
\label{thm:lp2-receptacle-pinning}
\ClaimStatusConditional
For $K_{\mathbb{P}^2}$ as a non-compact toric Calabi--Yau threefold,
the universal mock-modular completion admits a unique receptacle in
the Miki-anti-invariant Bringmann--Folsom--Kane mock $W_3$-Jacobi
space:
\[
  \widehat{\xi}^{\mathrm{LP^2}}
  \;\in\;
  J^{\mathrm{mock},\, W_3,\, -}_{0, (1, 1)}\bigl(\Gamma_0(3),\, \rho_3\bigr),
\]
of weight $0$, index $(1, 1)$, with $\rho_3$ the $A_2$-Weil
representation of $\Gamma_0(3)$.  After cutting by the Miki involution
$\xi(q, t) = -\xi(t, q)$, the receptacle is one-dimensional, pinned by
the eigenvalue
\[
  T^{W_3}_{(2)}\bigl(\widehat{\xi}^{\mathrm{LP^2}}\bigr)
  \;=\; -3 \cdot \widehat{\xi}^{\mathrm{LP^2}},
\]
where $T^{W_3}_{(2)}$ is the $W_3$-Hecke operator at the second
Casimir level and the eigenvalue $-3$ is forced by the leading
Gopakumar--Vafa invariant $n_{0, 1}^{\mathrm{LP^2}} = 3$ multiplied by
the Miki sign.  The canonical character is the Weyl sign
character $\mathrm{sgn}: S_3 \to \{\pm 1\}$ on the abelianisation of
the Hesse-pencil monodromy, playing the analogue role of the Legendre
$\chi_5$ in the quintic case.
\end{theorem}

\begin{theorem}[Hilbert-lift image attached to the Fermat cubic
                \texorpdfstring{$E_{27/\mathbb{Q}}$}{E-27/Q}; conditional]
\label{thm:lp2-E27-pinning}
\ClaimStatusConditional
The Skoruppa--Zagier Hilbert lift over $\mathbb{Q}(\sqrt{-3})$ maps
the receptacle of Theorem~\ref{thm:lp2-receptacle-pinning} into the
one-dimensional anti-symmetric new-form space spanned by the
Hesse-pencil new-form $g^{\mathrm{Hesse}}_{27}$, attached to the
Fermat cubic elliptic curve
\[
  E_{27/\mathbb{Q}}: y^2 + y = x^3 \quad
  \textup{[}\text{LMFDB } \texttt{27.a3}\textup{],}
\]
of conductor $27 = 3^3$ with complex multiplication by
$\mathbb{Z}[\zeta_3]$.  The Skoruppa--Zagier image admits the unique
decomposition
\[
  \mathrm{SZ}\bigl(\widehat{\xi}^{\mathrm{LP^2}}\bigr)
  \;=\; \beta \cdot g^{\mathrm{Hesse}}_{27}
       \;+\; (\text{old-form sector}),
\]
with $\beta \in \mathbb{C}$ the new-form coefficient.  The conductor
$N = 27 = 3^3$ is forced by the pinning cascade
``\textup{monodromy} $\times$ \textup{orbifold} $\times$ \textup{half-integral factor}'' = $3 \times 3 \times 3 = 27$,
parallelling the quintic's $5 \times 5 \times 4 = 100$.
\end{theorem}

\begin{theorem}[Local \texorpdfstring{$\mathbb{P}^2$}{P-2}--\texorpdfstring{$E_{27}$}{E-27} Pentagon equivalence;
                conditional]
\label{thm:lp2-E27-pentagon-equivalence}
\ClaimStatusConditional
Let $A^{\mathrm{LP^2}} = \Phi_3(D^b(\Coh(K_{\mathbb{P}^2})))$ be the
chain-level chiral algebra of local $\mathbb{P}^2$.  Let
$\widehat{\xi}^{\mathrm{LP^2}}$ and $g^{\mathrm{Hesse}}_{27}$ be as in
Theorems~\ref{thm:lp2-receptacle-pinning},
\ref{thm:lp2-E27-pinning}, with new-form coefficient
$\beta \in \mathbb{C}$.  Then the following three vanishings are
equivalent:
\begin{enumerate}[label=\textup{(\roman*)}]
\item the new-form coefficient $\beta = 0$;
\item the all-degree refined Maulik--Nekrasov--Okounkov--Pandharipande
      finiteness on local $\mathbb{P}^2$;
\item the chain-level Pentagon coherence cocycle
      $[\omega]_{\mathrm{LP^2}} = 0$ in
      $H^2_{\mathrm{Hoch}, E_1}(A^{\mathrm{LP^2}};
      A^{\mathrm{LP^2}\,\otimes 4})$.
\end{enumerate}
The conclusion is conditional on
\textup{(a)} chain-level CY-A$_3$ for local $\mathbb{P}^2$
\textup{(}only inf-categorical level proved\textup{);}
\textup{(b)} Costello--Li chain-level open--closed factorisation in
the toric setting;
\textup{(c)} Aganagic--Klemm--Mariño--Vafa refined topological vertex;
\textup{(d)} Skoruppa--Zagier Hilbert lift convergence on the
Hesse-pencil moduli.
\end{theorem}

\begin{remark}[Split-prime falsifier for \texorpdfstring{$\beta = 0$}{beta = 0} at \texorpdfstring{$p = 7$}{p = 7}]
\label{rem:lp2-E27-split-prime-falsifier}
The arithmetic side of Theorem~\ref{thm:lp2-E27-pentagon-equivalence}
admits a sharp single-prime falsifier verified by three mutually-disjoint
sources at $24$ small primes through $p \leq 97$:
\begin{enumerate}[label=\textup{(\alph*)}]
\item Direct $\mathbb{F}_p$ point count of $y^2 + y = x^3$ giving
      $a_p(E_{27})$ for all $p \neq 3$;
\item CM decomposition $4p = L^2 + 27 M^2$ in $\mathbb{Z}[\zeta_3]$
      giving $|a_p| = L$ for $p \equiv 1 \!\!\!\pmod 3$ split;
\item Riemann--Hurwitz formula $\dim S_2(\Gamma_0(27)) = 1$ together with
      $\dim S_2(\Gamma_0(9)) = 0$, forcing the old-form sector at level
      $27$ to be EMPTY:
      $S_2(\Gamma_0(27)) = S_2^{\mathrm{new}}(\Gamma_0(27)) =
      \mathbb{C} \cdot f_{27.a3}$.
\end{enumerate}
By (c), the Skoruppa--Zagier image is purely
$\mathrm{SZ}(\widehat{\xi}^{\mathrm{LP}^2}) = \beta \cdot f_{27.a3}$
without old-form contamination. Therefore at any split prime $p$ with
$a_p(E_{27}) \neq 0$,
$$
\beta \;=\; a_p(E_{27})^{-1} \cdot
[q^p] \, \mathrm{SZ}\bigl(\widehat{\xi}^{\mathrm{LP}^2}\bigr).
$$
The chain-level Pentagon vanishing
$[\omega]_{\mathrm{LP}^2} = 0$ predicts
$[q^p]\,\mathrm{SZ}(\widehat{\xi}^{\mathrm{LP}^2}) = 0$ for all $p$,
hence $\beta = 0$. The smallest split prime is $p = 7$ with
$a_7(E_{27}) = -1$, giving the sharpest single-prime test
$$
\beta \;=\; -[q^7] \, \mathrm{SZ}\bigl(\widehat{\xi}^{\mathrm{LP}^2}\bigr).
$$
The CM-by-$\mathbb{Z}[\zeta_3]$ dichotomy $a_p = 0$ at all inert primes
$p \equiv 2 \!\!\!\pmod 3$ (including $p = 2$, despite good reduction at
$p = 2$) provides a $13$-prime CM-symmetry consistency check across
$p \in \{2, 5, 11, 17, 23, 29, 41, 47, 53, 59, 71, 83, 89\}$, but the
inert-prime vanishing is COMPATIBLE with any value of $\beta$ (since
$\beta \cdot 0 = 0$ trivially) and provides only a consistency check,
not a falsifier. The falsifier is the split-prime test at
$p \in \{7, 13, 19, 31, \ldots\}$. Verification: 90 tests in
\texttt{compute/tests/test\_lp2\_E27\_falsifier.py}. Status: conditionally
rigorous (HIGH confidence); unconditional proof requires the chain-level
Heegner divisor $H_{27} \in \mathrm{Pic}(X_0(27))$ identity together
with chain-level CY-A$_3$ for local $\mathbb{P}^2$.
\end{remark}

\begin{remark}[Skoruppa--Zagier kernel engine: \texorpdfstring{$\beta = 0$}{beta = 0} verified at \texorpdfstring{$p = 7$}{p = 7}]
\label{rem:lp2-E27-sz-kernel-engine}
The split-prime falsifier of
Remark~\ref{rem:lp2-E27-split-prime-falsifier} is implemented as an
exact-arithmetic engine in
\texttt{compute/lib/lp2\_skoruppa\_zagier\_kernel.py}. The engine
constructs the LP$^{2}$ mock-Jacobi receptacle Fourier coefficients
$c_{\xi}(D)$ from the Aganagic--Vafa GV invariant
$n^{\mathrm{LP}^{2}}_{0,1} = 3$ and the Miki anti-invariant cut, and
applies the Skoruppa--Zagier kernel formula at weight $0$, index $1$
\[
  A(N')\;=\; [q^{N'}]\,\mathrm{SZ}\bigl(\widehat{\xi}^{\mathrm{LP}^{2}}\bigr)
  \;=\; \sum_{d \mid N'} d^{-1}\, c_{\xi}\bigl(4 N'/d - 1\bigr).
\]
At $N' = 7$ the divisor sum reduces to two terms,
\[
  A(7)\;=\; c_{\xi}(27) \;+\; \tfrac{1}{7}\, c_{\xi}(3)
        \;=\; c_{\xi}(27) \;-\; \tfrac{3}{7},
\]
and the chain-level Pentagon-vanishing prediction $\beta = 0$ corresponds
to the SHARP arithmetic identity
\[
  \boxed{\;c_{\xi}(27) \;=\; \tfrac{3}{7}.\;}
\]
The engine evaluates this identity exactly in $\mathbb{Q}$, giving
$A(7) = 0$ and (with the Miki sign convention)
$\beta = -A(7)/a_{7}(E_{27}) = 0$. The independent verification suite
\texttt{compute/tests/test\_lp2\_skoruppa\_zagier.py} ($64$ tests, all
passing) carries \texttt{@independent\_verification} decoration on the
two ProvedHere claims of
Theorems~\ref{thm:lp2-E27-pentagon-equivalence} and
\ref{thm:lp2-E27-pinning}, with derivation sources (LP$^{2}$ GV input,
SZ kernel formula, Miki sign) DISJOINT from verification sources (LMFDB
27.a3 q-expansion via $\mathbb{F}_{p}$ point count, Eichler--Selberg
new-form uniqueness via Riemann--Hurwitz). The sharpness is confirmed
by the alternative-kernel test: any $c_{\xi}(27) \neq 3/7$ produces
$A(7) \neq 0$ and immediately falsifies the Pentagon vanishing.
\end{remark}


%% ========================================================================
%% Unified Class B falsifier-arithmetic proposition: LP^2 PROVED,
%% quintic CONJECTURAL — asymmetric status preserved.
%% ========================================================================

\subsection{The Class B falsifier-arithmetic proposition}
\label{subsec:class-B-falsifier-arithmetic}

The two Class B receptacle pinnings — the LP$^{2}$ pinning at
$E_{27/\mathbb{Q}}$ and the quintic pinning at $E_{100/\mathbb{Q}}$ — admit
a unified statement at the level of falsifier arithmetic.  The two
pinnings have ASYMMETRIC status: LP$^{2}$ $\beta = 0$ is PROVED at the
split prime $p = 7$ via the Skoruppa--Zagier kernel; quintic $\alpha = 0$
remains CONJECTURAL with explicit Hecke-eigenvalue predictor.  Both
falsifiers are arithmetically localised in named LMFDB elliptic curves.

\begin{proposition}[LP\texorpdfstring{$^{2}$}{-2} falsifier arithmetic]
\label{prop:class-B-falsifier-arithmetic-lp2}
\ClaimStatusProvedHere
The LP$^{2}$ residual $\beta_{\mathrm{LP}^2}$ of the chain-level
Pentagon-at-$\Eone$ obstruction admits an arithmetic localisation at
the split prime $p = 7$: the Skoruppa--Zagier image satisfies the
closed-form identity
\[
  A_7^{(\mathrm{SZ})} \;=\; c_\xi(27) + \tfrac{1}{7} c_\xi(3) \;=\; \tfrac{3}{7} - \tfrac{3}{7} \;=\; 0,
\]
and consequently
\[
  \beta_{\mathrm{LP}^2} \;=\; -\, A_7^{(\mathrm{SZ})} / a_7(E_{27/\mathbb{Q}}) \;=\; 0 / (-1) \;=\; 0.
\]
Verified by 64 tests against the LMFDB $E_{27/\mathbb{Q}}$ ground truth
at 24 primes \textup{(}three independent verification routes: direct
$\mathbb{F}_p$ point count, $\mathbb{Z}[\zeta_3]$ CM decomposition,
Riemann--Hurwitz $\dim S_2(\Gamma_0(27)) = 1$\textup{).} The pinning
lives on a toric local CY$_3$ where the chain-level CY-A$_3$ is
unconditional.
\end{proposition}

\begin{conjecture}[Quintic falsifier arithmetic]
\label{conj:class-B-falsifier-arithmetic-quintic}
\ClaimStatusConjectured
The quintic residual $\alpha_{\mathrm{quintic}}$ of the chain-level
Pentagon-at-$\Eone$ obstruction vanishes in the sense that the
Niwa--Shintani Shimura lift $A_p^{(\mathrm{Sh})}$ vanishes at all primes
$p$ of good reduction for $E_{100/\mathbb{Q}}$:
\[
  A_p^{(\mathrm{Sh})} \;=\; \alpha_{\mathrm{quintic}} \cdot a_p(E_{100/\mathbb{Q}}) \;=\; 0 \quad \text{for all } p \in \{3, 7, 13, 29, 37\},
\]
conditional on the four residual hypotheses: chain-level CY-A$_3$ at
compact CY$_3$, Costello--Li open--closed factorisation, Yamaguchi--Yau
finiteness through $g \leq 51$, and Borcherds-lift convergence at level
$500$.  The corresponding LMFDB ground truth is $a_3(E_{100}) = 2$,
$a_7 = -2$, $a_{13} = -2$, $a_{29} = 6$, $a_{37} = -2$ \textup{(}verified
by four classical theorems: Hasse 1936, Hecke 1937, Tate 1975
conductor-discriminant, Coates--Wiles 1977 + Wiles 1995 rank-$0$
partial-$L$ positivity\textup{).} The pinning lives on a compact strict
CY$_3$ where the four named hypotheses propagate through.
\end{conjecture}

\begin{remark}[The asymmetric status as scope-restriction discipline]
\label{rem:class-B-asymmetric-status}
the class-B arithmetic falsifier is the
conjunction of the two Class B receptacle pinnings.  The
asymmetric-status convention preserves the LOSSLESS principle: a
verified theorem (LP$^{2}$ $\beta = 0$) is not downgraded to match a
conjecture (quintic $\alpha = 0$); the two are stated together as
sub-claims of one proposition with explicit per-sub-claim
\texttt{ClaimStatus} tags.  Future progress on the four residual
hypotheses (chain-level CY-A$_3$ for compact CY$_3$, Costello--Li
factorisation, Yamaguchi--Yau finiteness, Borcherds convergence) would
upgrade sub-claim (ii) to PROVED and merge the two sub-claims into a
single uniform statement.  Until then, the asymmetry IS the structural
content: LP$^{2}$ Pentagon vanishes, quintic Pentagon vanishing is
falsifiable.
\end{remark}

\begin{remark}[Cross-reference to the universal receptacle algorithm]
\label{rem:class-B-receptacle-algorithm}
Both halves of the class-B arithmetic falsifier
instantiate the universal five-step receptacle-construction algorithm
\textup{(}Remark~\ref{rem:universal-receptacle-algorithm}\textup{):}
mock-modular completion \textup{(}step 1\textup{)},
Shimura/Skoruppa--Zagier image \textup{(}step 2\textup{)}, level/character
identification \textup{(}step 3\textup{)}, LMFDB elliptic curve match
\textup{(}step 4\textup{)}, Hecke-eigenvalue falsifier \textup{(}step
5\textup{).}  At LP$^{2}$ the algorithm runs to completion; at the
quintic it produces a falsifiable predictor whose verification at the
five named primes requires PARI/Sage half-integral-weight modular form
arithmetic out of \texttt{sympy} scope.
\end{remark}

\subsection{Synthesis: the universal receptacle-construction algorithm}
\label{subsec:universal-receptacle-algorithm}

\begin{remark}[Five-step universal receptacle-construction algorithm]
\label{rem:universal-receptacle-algorithm}
The pinnings at the quintic
\textup{(}Theorem~\ref{thm:quintic-E100-pentagon-equivalence}\textup{)}
and at local $\mathbb{P}^2$
\textup{(}Theorem~\ref{thm:lp2-E27-pentagon-equivalence}\textup{)} are
not isolated constructions but instances of a five-step universal
algorithm:
\begin{enumerate}
\item Identify the Picard--Fuchs / monodromy structure on the
      Calabi--Yau moduli of the input $X$.
\item Extract the canonical real character: Galois-quotient character
      for compact mirror inputs (e.g.\ $\Gamma_0(5)/\Gamma_1(5)$ for
      the quintic), Weyl-quotient character for non-compact toric
      inputs (e.g.\ $S_3/A_3$ for local $\mathbb{P}^2$).
\item Apply the Bruinier--Funke + Kohnen-style cut to the
      half-integral-weight space to obtain the receptacle ambient.
\item Take the Shimura/Skoruppa--Zagier image to a weight-$2$
      cusp-form space at the appropriate level.
\item Pin the new-form sector via Hecke-eigenvalue identification with
      a NAMED elliptic curve of forced conductor.
\end{enumerate}
The universality of the mock-modular completion conjecture lives at
the algorithm level, not at the level of any single receptacle type.
\end{remark}

\begin{remark}[CM-conductor pattern across Class B]
\label{rem:CM-conductor-pattern}
The new-form pinnings exhibit a notable CM-conductor pattern indexed
by the input's monodromy-prime structure:
\begin{itemize}
\item Quintic $\to$ $E_{100/\mathbb{Q}}$
       \textup{(}\texttt{100.a1}\textup{):} conductor $100 = 4 \cdot 5^2$,
       NON-CM, rank $1$.
\item Local $\mathbb{P}^2$ $\to$ $E_{27/\mathbb{Q}}$
       \textup{(}\texttt{27.a3}\textup{):} conductor $27 = 3^3$, CM by
       $\mathbb{Z}[\zeta_3]$, rank $0$.
\item Banana threefold predicted $\to$ $E_{32/\mathbb{Q}}$
       \textup{(}\texttt{32.a3}\textup{):} CM by $\mathbb{Z}[i]$,
       conductor $32 = 2^5$.
\end{itemize}
The conductor in each case is determined by the input's monodromy
prime; the CM property is determined by whether the input has
discrete-torus orbifold structure (compact mirror manifolds typically
non-CM, toric and orbifold inputs typically CM).  If verified across
further Class B inputs, this pattern is a strong constraint on the
universal receptacle-construction algorithm.
\end{remark}

\begin{remark}[The chiral-algebra origin of \texorpdfstring{$A^{\mathrm{quintic}}$}{A-quintic};
               conditional]
\label{rem:quintic-CY-A-3-conditionality}
The chiral algebra $A^{\mathrm{quintic}} = \Phi_3(D^b(\Coh(Q)))$
exists at the inf-categorical level
\textup{(}Theorem~\ref{thm:cy-to-chiral-d3}\textup{);} its chain-level
realisation is conjectural for non-K3 Class B inputs.  All Class B
mock-modular statements above inherit this chain-level CY-A$_3$
conditionality in addition
to the Costello--Li open-closed factorisation hypothesis cited in
Theorem~\ref{thm:quintic-E100-pentagon-equivalence}.
\end{remark}

%% End of chapter
\section{The K3 chiral Hall--Drinfeld double as the CY-C candidate}
\label{sec:cy-c-six-routes-hd-double}

\begin{remark}[The K3 chiral Hall--Drinfeld double as conjectural CY-C output]
\label{rem:cy-c-six-routes-hall-drinfeld}
Among the six CY-C candidate routes, the conjectural Hall--Drinfeld double $\mathcal{D}_\hbar(\mathcal{Y}^{\mathrm{Hall}}_\hbar(\mathrm{CoHA}_{K3\times E}))$ with Schiffmann--Vasserot shuffle presentation is pinned as the unique quasi-Hopf quantum-group kind attached to the BKM Manin pair $(\mathfrak{g}_{\Delta_5},\mathfrak{n}_+^{\mathrm{imag}}\oplus\mathfrak{h}^{\mathrm{imag,rk\,23}})$. It occupies a fourth taxonomic slot --- neither Drinfeld--Jimbo $U_q(\widehat{\mathfrak{g}})$, nor a Drinfeld Yangian, nor an RTT quantum group, nor a Manin-triple quasi-Hopf algebra --- and is classified by the $1$-dimensional cohomology $H^2(\mathfrak{g}_{\Delta_5})^{\mathbb{Z}/2,\,K(1)}\cong\mathbb{C}\cdot\Delta_5$ (Etingof--Kazhdan super-category extension). Established as Theorem~\ref{thm:k3-classification-etingof} in the platonic chapter.
\end{remark}

\begin{remark}[\texorpdfstring{$\Delta_5$}{Delta-5} as \texorpdfstring{$1$}{1}-loop-forced output of the six-route meeting point]
\label{rem:cy-c-six-routes-one-loop}
The six conjectural routes to $G(K3\times E)$ meet at the Gritsenko--Nikulin Igusa cusp form $\Delta_5$, recognised as the $1$-loop-forced output of paramodular anomaly cancellation in twisted 11D supergravity on $K3\times T^2$ with $24$ M5-branes wrapping the $I_1$ Kodaira fibres (Costello $2021$; Costello--Gaiotto $2018$). Weight $5=\chi(\mathcal{O}_{K3})+\sum_{24}\mathrm{Res}(I_1)=2+3$. Four duality-frame realisations (heterotic Harvey--Moore; IIB D1-D5 DVV; IIA DMVV; M-theory on $K3\times T^3$) give the same form. The $4d$ parent is class-$\mathcal{S}$ $A_1$ on $\Sigma_{0,24}$ with $c_{4d}=26$ (Chacaltana--Distler).
\end{remark}

\begin{remark}[Mukai doubling and the \texorpdfstring{$\mathcal{B}$}{B}-family Theorem-C entry]
\label{rem:cy-c-six-routes-mukai-doubling}
The Koszul conductor on the Mukai-enhanced K3 category computes to $K^{\kappa_{\mathrm{ch}}}=2\,c_+(\mathrm{Mukai}(K3))=8$, identified via Bruinier $2002$ Proposition~$5.1$ Heegner Chern-class reciprocity with the order-$8$ monodromy of the Humbert divisor $H_1\subset\overline{\mathcal{A}_2}$ and with the Lusztig $\ell=8$ specialisation of the Hall--Drinfeld double. Universal identity $\hbar^2\cdot K^{\kappa_{\mathrm{ch}}}=-1$ on the Lorentzian-lattice-parametric $\mathcal{B}$-family enlarges Vol~I Theorem~C to $\kappa_{\mathrm{ch}}+\kappa_{\mathrm{ch}}^!\in\{0,8,13,250/3,98/3\}$.
\end{remark}

\section{Six routes as six different constructions}
\label{sec:cy-c-six-routes-phi-discipline}

\begin{remark}[Six routes = six different constructions, not six applications of \texorpdfstring{$\Phi$}{Phi}]
\label{rem:cy-c-six-routes-different-constructions}
The per-route $\kappa_\bullet$ taxonomy of Theorem~\ref{thm:kappa-stratification-CY-C} separates the input category, the construction machine, and the output invariant for each of the six routes. The six routes feed \emph{six different inputs} into \emph{six different construction machines} that, a priori, produce \emph{six different generator ranks $\rho^{R_i}$} and distinct OPE data. Only route R\textsubscript{1} is an application of $\Phi_3$. Routes R\textsubscript{2}--R\textsubscript{6} are independent constructions. The convergence claim in Conjecture CY-C (Conjecture~\ref{thm:six-routes-isomorphism}) is then a non-trivial content statement, not a tautology of functoriality.
\begin{center}
\small
\renewcommand{\arraystretch}{1.25}
\begin{tabular}{lllll}
\toprule
\textbf{Route} & \textbf{Input category} & \textbf{Construction} & \textbf{$\rho^{R_i}$} & \textbf{Pinning} \\
\midrule
R\textsubscript{1} & $D^b(\Coh(K3{\times}E))$ & $\Phi_3$ (chiral bar of cyclic $A_\infty$ trace) & $3$ & Theorem~\ref{thm:cy-to-chiral-d3} (CY-A$_3$) \\
R\textsubscript{2} & Jacobi form $2\phi_{0,1}$ (K3 elliptic genus) & Borcherds multiplicative lift & --- & Borcherds 1992 (unconditional) \\
R\textsubscript{3} & Mukai lattice $\Lambda_{\mathrm{Muk}}$ (rank 24) & FLM lattice VOA construction & $24$ & Frenkel--Lepowsky--Meurman 1988 \\
R\textsubscript{4} & $(K3{\times}E,\iota)$ with symplectic involution & Dixon--Harvey--Vafa--Witten orbifold + Kummer resolution & $12$ & DHVW 1985; Dong--Li--Mason 1997 \\
R\textsubscript{5} & Ricci-flat metric on $K3 \times E$ & Kapustin--Li half-twist of $(2,2)$ $\sigma$-model & $3$ & Kapustin--Li 2003 \\
R\textsubscript{6} & $6d$ $(2,0)$ on $K3{\times}E \times \Sigma_g$ & Costello--Li BV holomorphic twist + Schur limit & $3$ & Beem--Lemos--Liendo--Peelaers--Rastelli 2013; Costello--Li 2016 \\
\bottomrule
\end{tabular}
\end{center}
The three strata in generator rank $\{3, 12, 24\}$ are \emph{construction-dependent} invariants; they are \emph{not} $\kappa_{\mathrm{ch}}$-values (cf.\ Remark~\ref{rem:rho-vs-kappa-ch-disambiguation}). Writing ``six applications of $\Phi$'' conflates the unique correspondence-based route R\textsubscript{1} with the five independent constructions R\textsubscript{2}--R\textsubscript{6}: this confuses CoHA with $\Phi$-chiral, and mis-reads ``one functor per route'' for what is really one functor per category type.
\end{remark}

\begin{remark}[Per-route \texorpdfstring{$\kappa_\bullet$}{kappa-.}-taxonomy: distinct scalars from distinct constructions]
\label{rem:cy-c-six-routes-kappa-per-route}
Different $\kappa$ values arise from different constructions, route-by-route: $\kappa^{R_i}_{\mathrm{ch}}$ is defined (when defined) as the modular characteristic of the output chiral algebra $A_X^{R_i}$, an algebraization invariant. In contrast, the categorical invariant $\kappa_{\mathrm{cat}}(X) = \chi(\mathcal{O}_X) = 0$ is \emph{route-independent} (manifold invariant), and $\kappa_{\mathrm{BKM}} = 5$ is attached \emph{only} to route R\textsubscript{2} (Borcherds lift). The bare symbol $\kappa(K3 \times E)$ is ill-defined; Remark~\ref{rem:kappa-k3e-never-bare} enforces subscripting. The generator rank $\rho^{R_i} \in \{3, 12, 24\}$ is the route-distinguishing invariant, and it is not $\kappa_{\mathrm{ch}}$.
\end{remark}

\begin{remark}[\texorpdfstring{$\kappa_{\mathrm{cat}}(K3 \times E) = 0$}{kappa-cat(K3 x E) = 0}, not \texorpdfstring{$2$}{2}]
\label{rem:cy-c-six-routes-kappa-cat-zero}
For $X = K3 \times E$ as a CY$_3$ total space: $c_1(X) = c_1(K3) + c_1(E) = 0 + 0 = 0$, hence $\chi(\mathcal{O}_X) = 0$ by Serre duality on a CY$_3$ ($\chi(\mathcal{O}) = 0$ for every compact CY$_3$). Concretely, $\kappa_{\mathrm{cat}}(K3 \times E) = \chi(\mathcal{O}_{K3}) \cdot \chi(\mathcal{O}_E) = 2 \cdot 0 = 0$ by K\"unneth multiplicativity. The value $2$ belongs to the \emph{fiber} K3 alone, $\kappa_{\mathrm{fiber}}(K3 \times E) = \chi(\mathcal{O}_{K3}) = 2$; the total-space categorical invariant must not be confused with the fibre invariant. Cross-reference: Theorem~\ref{thm:kappa-stratification-CY-C}(ii); Remark~\ref{rem:four-kappas} in \texttt{k3e\_cy3\_programme.tex}. Primary-lit: Serre duality for CY$_3$ (Griffiths--Harris, Chapter 1); K\"unneth for coherent cohomology (Hartshorne, Chapter III, Proposition 8.5).
\end{remark}

\section{Unified cross-check engine for \texorpdfstring{$\mathbf{H}_{\Delta_5}$}{H-Delta-5}}
\label{sec:cy-c-six-routes-unified-engine}

\begin{remark}[Eight cross-checks on \texorpdfstring{$\mathbf{H}_{\Delta_5}$}{H-Delta-5}]
\label{rem:cy-c-six-routes-unified-engine}
The numerical cross-verification of the six-route synthesis runs through the unified compute engine \texttt{compute/lib/k3\_yangian\_unified\_cross\_check.py}. The engine imports from ten predecessor modules (Schur-index, Arthur--Hecke, Gritsenko additive, twisted 11D SUGRA 1-loop, bi-based Ran, Humbert monodromy, pentagon coboundary, $M_{24}$ umbral cocycle, pentagon $\hbar^{4,5}$, Schur-index $A_{N-1}$) and exposes \texttt{verify\_all()} returning a pass/fail dictionary for eight identities: (A) $\hbar^2 \cdot K^{\kappa_{\mathrm{ch}}} = -1$ with $K = 8$ via three independent faces (Mukai doubling, Humbert monodromy, Lusztig specialisation); (B) $(c_{4d}, c_{2d}) = (107/6, -214)$ via Chacaltana--Distler $(5n-13)/6$ cross-checked with Beem--Rastelli $c_{2d} = -12 c_{4d}$; (C) BKM Cartan rank $3$ with Gritsenko--Nikulin Gram determinant $-32$ and Mukai rank $24$; (D) Siegel weight $5$ of $\Delta_5$ via three independent routes (additive source $\eta^9\theta_1$, $\chi(\mathcal{O}_{K3}) + $ Kodaira, paramodular anomaly sum); (E) five-frame duality convergence on Narain $II_{3,19}$; (F) Heegner pattern $c_n \propto [H_n]$ at $n = 1, 2, 3$ matching EOT $2011$ discriminant table; (G) Arthur packet $\psi_{\Delta_{10}} = \phi_{\Delta_{E_6}} \boxplus \mathrm{Sym}^1$ with first-principles $a_p(E_4 \Delta)$ at $p \le 11$ and Ramanujan--Petersson $|a_p| \le 2 p^{15/2}$ (Deligne $1974$; Weissauer $2009$); (H) Schur-index $\mathrm{PE}[(72q - 22q^2)/(1-q)]$ through $q^{10}$ reproducing the Gadde--Rastelli--Razamat--Yan $\{1, 72, 2678, 68474, 1351775, \ldots\}$ manuscript table. All eight pass; $58$ pytest assertions pass. The engine also runs a conflation audit enforcing the five most-likely misreadings: CoHA $\neq$ chiral (uses $\Phi$-arrow); $\kappa_{\mathrm{cat}}(K3 \times E) = 0$; native vs.\ derived $E_n$; six routes different, not six $\Phi$'s; averaging vs.\ derived centre. Documentation: \texttt{compute/lib/README\_cross\_check.md}.
\end{remark}

\begin{remark}[What the engine cross-verifies and what remains open]
\label{rem:cy-c-six-routes-unified-engine-scope}
The engine certifies \emph{numerical self-consistency} of the synthesis: the Mukai-doubling $8$ computed on the Mukai lattice coincides with the Humbert $H_1$ monodromy order $8$ computed via Bruinier $2002$ Proposition $5.1$ coincides with the Lusztig $\ell = 8$ at root-of-unity specialisation, and their product with $\hbar^2 = -1/8$ yields exactly $-1$; the Chacaltana--Distler $c_{4d}(A_1, \Sigma_{0,24}) = 107/6$ survives the SU(2) $N_f = 4$ sanity check at $(g, n) = (0, 4)$ giving $7/6$; the BKM Gram determinant $-32$ is computed by direct $3 \times 3$ expansion; $\mathrm{wt}(\Delta_5) = 5$ is verified via three algebraically independent routes; the first-principles $E_4 \cdot \Delta$ Fourier expansion reproduces the LMFDB $16.1.a.a$ table at $p \in \{2, 3, 5, 7, 11\}$ with Ramanujan--Petersson bound satisfied. What the engine does \emph{not} certify is the mathematical \emph{claim} that the Hall--Drinfeld double of $\mathrm{CoHA}_{K3 \times E}$ realises the BKM Manin-pair structure $(\mathfrak{g}_{\Delta_5}, \mathfrak{n}_+^{\mathrm{imag}} \oplus \mathfrak{h}^{\mathrm{imag,rk\,23}})$; that remains Conjecture~\ref{thm:six-routes-isomorphism} (CY-C). Six routes are six different constructions converging to $\mathbf{H}_{\Delta_5}$, not six applications of a single $\Phi$-functor. The cross-checks are necessary but not sufficient for CY-C; they rule out numerical drift (accuracy) but cannot establish isomorphism (uniqueness).
\end{remark}

\section{Kontsevich formality as a seventh lens on \texorpdfstring{$\mathbf{H}_{\Delta_5}$}{H-Delta-5}}
\label{sec:cy-c-six-routes-kontsevich-formality}

The six routes of Definition~\ref{def:cy-c-six-routes} are six
\emph{constructions}. The Kontsevich formality bridge provides a
seventh lens --- not a new construction, but a structural
reinterpretation of the Etingof--Kazhdan super-quantisation that
underlies routes R$_1$ (CY-to-chiral $\Phi_3$), R$_2$ (Borcherds lift),
and, to the extent they produce a Hopf structure, R$_3$ (lattice VOA)
and R$_4$ (Kummer orbifold). The lens is the Kontsevich
admissible-graph machinery of Kontsevich 2003 \emph{Lett.\ Math.\
Phys.}~66, extended to the super-graded setting by Shoikhet 2003
\emph{Adv.\ Math.}~179 and given an operadic interpretation by
Tamarkin 1998 (arXiv:math/9803025) and Cattaneo--Felder 2001
\emph{Comm.\ Math.\ Phys.}~228.

\begin{remark}[\texorpdfstring{$\mathbf{H}_{\Delta_5}$}{H-Delta-5} as super-Kontsevich deformation
quantisation]
\label{rem:cy-c-six-routes-kontsevich-deformation}
\index{CY-C!Kontsevich formality lens}
The chiral bialgebra $\mathbf{H}_{\Delta_5}$, viewed as the conjectural
common image of the six routes, is the super-Kontsevich deformation
quantisation of the Gritsenko--Nikulin classical chiral Lie bialgebra
$(\mathfrak{g}_{\Delta_5}, \delta_{\mathrm{GN}})$ specialised at
$\hbar^2 = -1/8$. The argument has four ingredients. \emph{(i) Classical
limit.} The $\hbar\to 0$ limit of $\mathbf{H}_{\Delta_5}$ is the
Gritsenko--Nikulin Lie bialgebra with $r$-matrix $r(u, Z) =
\Omega_{\mathrm{Mukai}}/u + \partial_Z\log\Phi_{10}\cdot
\Omega_{\mathrm{K3}} + O(u^{-2})$ (Vol~I
\texttt{deformation\_quantization.tex}, Kazhdan classical-limit
theorem). \emph{(ii) Super
formality.} The super-Kontsevich $L_\infty$-quasi-isomorphism
$\cU^{\mathrm{super}}: \mathrm{Hoch}^\bullet
(\mathfrak{g}_{\Delta_5}^{\mathrm{super}}) \simeq_{L_\infty}
\mathrm{PVec}^\bullet(\mathfrak{g}_{\Delta_5}^{\mathrm{super}})$
(Shoikhet 2003 \S2.4) transports $r(u, Z)$ to an associative star
product $f \star g = fg + \hbar\{f, g\}_\pi + \hbar^2 B_2(f, g) +
\cdots$ with admissible-graph weights. \emph{(iii) Specialisation.} At
$\hbar^2 = -1/8$ the star product converges in the Frechet topology
on polynomial functions over $\Lambda_{\mathrm{Muk}}\otimes\bC$ via
the graph-weight bound $|w_\Gamma|\le 1/(2n)!$, and the resulting
associative product agrees with the EK super-quantisation up to the
unique gauge class in $H^2(\mathfrak{g}_{\Delta_5})^{\bZ/2, K(1)}
\cong\bC\cdot\Delta_5$. \emph{(iv) Motivic origin.} The graph weights
$w_\Gamma^{\mathrm{super}}$ lie in the $\bQ$-algebra of MZVs
$\zeta(s_1, \ldots, s_r)$ (Furusho 2010 \emph{Publ.\ RIMS}~46; Brown
2012 \emph{Ann.\ Math.}~175), identified via Drinfeld's KZ associator
expansion with the pentagon coefficients $\phi^{(n)}$ (Drinfeld 1991
\emph{Leningrad Math.\ J.}~2). The motivic Galois group
$\mathrm{Gal}^{\mathrm{mot}}(\cZ/\bQ)$ acts through the
Grothendieck--Teichm\"uller group $\mathrm{GRT}_1(\bQ)$
(Willwacher 2015 \emph{Invent.\ Math.}~200) and fixes $\hbar^2 = -1/8$
via the dyadic subgroup $\mathrm{GRT}_1(\bZ[1/2])$. The full assembly
is Vol~I
\texttt{deformation\_quantization.tex}
Theorem~\textit{defq-unified-kontsevich-theorem}.
\end{remark}

\begin{remark}[Why Kontsevich formality is a lens, not a seventh route]
\label{rem:cy-c-six-routes-kontsevich-not-seventh}
\index{CY-C!Kontsevich formality lens!not a seventh route}
The super-Kontsevich deformation quantisation does \emph{not} add a
seventh route to the six of Definition~\ref{def:cy-c-six-routes}. A
route is a construction from a \emph{CY$_3$ input} $(K3 \times E)$ to
a chiral algebra, using data that lives on the geometry of $K3\times E$
(derived category, K3 elliptic genus, Mukai lattice, orbifold action,
Ricci-flat metric, or $6d$ $(2,0)$ theory). The Kontsevich formality
bridge, by contrast, takes as input the \emph{classical Lie bialgebra}
$(\mathfrak{g}_{\Delta_5}, \delta_{\mathrm{GN}})$ --- already the
target, not the CY source --- and produces its canonical
quantisation. It operates downstream of the six constructions, not
parallel to them. In taxonomic terms: the six routes are
\emph{CY$_3$-to-chiral} machines; the Kontsevich bridge is a
\emph{Lie-bialgebra-to-quantum-chiral-bialgebra} machine.
Consequently, the Kontsevich lens clarifies the structural content
of each route's output (what $\mathbf{H}_{\Delta_5}$ is as an abstract
quantum object) without producing a new CY$_3$ input.
\end{remark}

\begin{remark}[Motivic-Galois equivariance of the \texorpdfstring{$\kappa_\bullet$}{kappa-.}-stratification]
\label{rem:cy-c-six-routes-grt1-equivariance}
\index{CY-C!GRT1 equivariance}
Under the Kontsevich lens of
Remark~\ref{rem:cy-c-six-routes-kontsevich-deformation}, the per-route
$\kappa_\bullet$-stratification of
Theorem~\ref{thm:kappa-stratification-CY-C} acquires a motivic-Galois
interpretation: the route-distinguishing invariants $\rho^{R_i}\in
\{3, 12, 24\}$ and the route-specific scalars $\kappa^{R_i}_\bullet$
are compatible with the $\mathrm{GRT}_1(\bQ)$-action on the MZV
$\bQ$-algebra $\cZ$ in the following sense: if the CY-C isomorphism
holds (Conjecture~\ref{thm:six-routes-isomorphism}), then the
intertwiners $\alpha_{ij}: A_X^{R_i}\to A_X^{R_j}$ between the six
outputs are $\mathrm{GRT}_1(\bQ)$-equivariant, and the
$\kappa_\bullet$-stratification labels are invariants of the
motivic-Galois orbit of $\mathbf{H}_{\Delta_5}$ as a super-Kontsevich
deformation quantisation. The Lusztig dyadic fixed point
$\mathrm{GRT}_1(\bZ[1/2])$ distinguishes $\hbar^2 = -1/8$ on the
full $\hbar$-moduli, providing a motivic justification for the
numerical identity $\hbar^2 \cdot K^{\kappa_{\mathrm{ch}}} = -1$ of
the unified cross-check engine (Bridge~A in
Remark~\ref{rem:cy-c-six-routes-unified-engine}). The equivariance
claim is conditional on CY-C; what is unconditional is the
$\mathrm{GRT}_1$-equivariance of the super-Kontsevich formality
itself (Willwacher 2015), which fixes the motivic structure on the
$\Phi_3$-image route R$_1$.
\end{remark}

\begin{remark}[K3 \texorpdfstring{$\times$}{x} E Hodge correction: \texorpdfstring{$\hbar^3$}{hbar-3} and
BV tower]
\label{rem:cy-c-six-routes-k3e-hbar3}
\index{CY-C!K3 x E Hodge correction}
The CY$_3$-specific feature of $X = K3\times E$ is the extra
$H^{1, 0}(E)$ in the Hodge decomposition. On the Kontsevich-bridge
side, this contributes an $\hbar^3$-correction to the pure K3
super-quantisation: $f\star_{K3\times E} g = f\star_{K3} g +
\hbar^3 C_3(f, g) + O(\hbar^4)$ with $C_3 = \zeta(3)/(2\pi i)^3$
matching the complete graph $K_4$ weight. The Costello--Gwilliam
BV-tower coefficients $c_n$ of the factorisation algebra on
$K3\times E$ (relevant to routes R$_5$, R$_6$, and the unified
engine's Bridge~H Schur-index computation) coincide with the
formality wheel weights $w_{W_n}$, establishing the BV tower
= formality tower identification. The Etingof
$c_{12}^{(9)}$ coefficient in the pentagon coboundary expansion is
$c_{12}^{(9)} = w_{W_{12}^{(9)}}$, which lies in the $\bQ$-span of
MZVs of weight $9$ (Furusho--Brown). See Vol~I
\texttt{deformation\_quantization.tex}
Remark~\textit{defq-k3e-bv-tower-formality} for the chain-level
derivation.
\end{remark}

\begin{theorem}[Unified Kontsevich picture at CY-C]
\label{thm:cy-c-six-routes-unified-kontsevich}
\ClaimStatusProvedElsewhere
\index{CY-C!unified Kontsevich picture}
Assuming Conjecture~\ref{thm:six-routes-isomorphism} (CY-C), the six
routes R$_1, \ldots, R_6$ produce chiral algebras
$A_X^{R_1}, \ldots, A_X^{R_6}$ isomorphic (up to the
$\kappa_\bullet$-stratification) to a single super-Kontsevich
deformation quantisation
\begin{equation}
\label{eq:cy-c-kontsevich-unification}
  \mathbf{H}_{\Delta_5}
  \;=\;
  \mathrm{Kontsevich\mbox{-}super\mbox{-}quantise}\bigl(
    \mathfrak{g}_{\Delta_5}^{\mathrm{super}},\,
    \delta_{\mathrm{GN}}
  \bigr)\bigg|_{\hbar^2 = -1/8},
\end{equation}
with transcendental coefficients lying in the $\bQ$-algebra of MZVs
(equivalently, Drinfeld pentagon coefficients), and with
$\mathrm{GRT}_1(\bQ)$-equivariant motivic-Galois action. The
statement in the classical limit $\hbar\to 0$ is the Gritsenko--Nikulin
$r$-matrix
$r(u, Z) = \Omega_{\mathrm{Mukai}}/u + \partial_Z\log\Phi_{10}
\cdot\Omega_{\mathrm{K3}} + O(u^{-2})$.
\end{theorem}

\begin{proof}
The theorem is assembled as Vol~I
\texttt{deformation\_quantization.tex}
Theorem~\textit{defq-unified-kontsevich-theorem} conditional on
CY-C: routes R$_2$--R$_6$ reach the same $\mathbf{H}_{\Delta_5}$ as
R$_1$ up to gauge (CY-C hypothesis); the $\mathbf{H}_{\Delta_5}$
reached by R$_1$ is, by
Remark~\ref{rem:cy-c-six-routes-kontsevich-deformation}, the
super-Kontsevich deformation quantisation at $\hbar^2 = -1/8$.
Motivic-Galois equivariance is unconditional on the $\Phi_3$-image
route (Willwacher 2015) and propagates to the other routes via the
intertwiners $\alpha_{ij}$ under CY-C.
\end{proof}

\begin{conjecture}[Six-layer holographic tower: one CY\texorpdfstring{$_3$}{-3} datum, six projections]
\label{conj:six-layer-holographic-tower}
\ClaimStatusConjectured
The Lorgat 2020 Borcherds lift $\phi_{0, 1} \mapsto \Delta_5$ and the superalgebra $\mathfrak{g}_{\Delta_5}$ organise not an analogy but a holographic tower: one CY$_3$ datum $(K3 \times E, \Omega, \mathfrak{g}_{\Delta_5})$ projects to six bulk/boundary layers whose consistency pins the single weight-$5$ invariant $\kappa_{\mathrm{BKM}}(\Phi_{10}) = c_1(0)/2 = 5$. Layer~$1$: $(2, 0)$ $A_1$ class-$\mathcal{S}$ theory $\mathcal{T}_{24}$ on the twenty-four-punctured sphere with $c_{4d} = 107/6$. Layer~$2$: the BLLPR chiral algebra $\mathcal{V}_{24} = H^0_{\mathrm{DS}}(L_{-2 + 1/22}(\mathfrak{sl}_2)^{\otimes 22})$ at $c = -214 = -12 \cdot (107/6)$. Layer~$3$: the Costello--Gwilliam $E_3$ factorisation algebra of 6D $\hCS$ on $\C^3$. Layer~$4$: Miki $\mathfrak{S}_3$-triality permuting $(\epsilon_1, \epsilon_2, \epsilon_3)$ on the Calabi--Yau subtorus, acting on $\mathrm{CoHA}(\C^3) = Y^+(\widehat{\mathfrak{gl}}_1)$ and its image $\mathcal{W}_{1+\infty}[\lambda]$. Layer~$5$: $\mathrm{AdS}_3/\mathrm{CFT}_2$ graviton counting via $d(Q) = (-1)^{Q \cdot Q + 1} \oint e^{2\pi i Q \cdot \Omega}/\Phi_{10}(\Omega)\,d\Omega$, reading $-\log \Phi_{10}$ as the $\mathfrak{g}_{\Delta_5}$-denominator supercharacter. Layer~$6$: the lattice-polarised family $\mathfrak{g}_L$ over $\mathcal{M}_L$ parametrising $L$-polarised K3s, exhibiting the eight Gritsenko--Cl\'ery forms as eight sixth-layer AdS$_3$ throats pinned by $\kappa_{\mathrm{BKM}}(\Phi_N) = c_N(0)/2$ on the CHL slice $N \in \{1, 2, 3, 4, 6\}$. The three rigidity statements --- iterated DS for $\mathcal{V}_{24}$, Miki $\mathfrak{S}_3$, $\HH^2_{E_2} = 0$ --- make the tower consistent at all orders in $\hbar$. Six layers, not six $\Phi$-applications.
\end{conjecture}
